%\input{tcilatex}


\documentclass{article}
\usepackage{amssymb}

%%%%%%%%%%%%%%%%%%%%%%%%%%%%%%%%%%%%%%%%%%%%%%%%%%%%%%%%%%%%%%%%%%%%%%%%%%%%%%%%%%%%%%%%%%%%%%%%%%%%
\usepackage{graphicx}
\usepackage{amsmath}

%TCIDATA{OutputFilter=LATEX.DLL}
%TCIDATA{Created=Tue Apr 29 19:13:34 2003}
%TCIDATA{LastRevised=Tue Jul 29 10:45:01 2003}
%TCIDATA{<META NAME="GraphicsSave" CONTENT="32">}
%TCIDATA{<META NAME="DocumentShell" CONTENT="General\Latex_article">}
%TCIDATA{CSTFile=article.cst}

\newcommand{\fig}[4]{
\begin{figure}
\begin{center}
\includegraphics[width = #4 cm]{#1}
\caption{#2}\label{#3}
\end{center}
\end{figure}
}
\addtolength{\voffset}{-2.5cm} \addtolength{\textheight}{5.0cm}
\addtolength{\hoffset}{-2.0cm} \addtolength{\textwidth}{5.0cm}
\input{tcilatex}

\begin{document}

\title{Bounce-Averaged Quasilinear Description of RFCD}
\date{April 29, 2003\\
\vspace{0.5cm} Updated July 2, 2003}
\author{Joan Decker and Yves Peysson}
\maketitle
\tableofcontents

\section{Original Formulation}

The quasilinear diffusion operator was developped by Kennel and Engelmann 
\cite{ken66}. The relativistic treatment was performed by Lerche \cite{ler68}%
, who proposes the following operator 
\begin{eqnarray}
Q\left( f\right) &=&-\frac{e^{2}}{\left( 2\pi \right) ^{3}}%
\lim_{V\rightarrow \infty }\frac{1}{V}\sum_{n=-\infty }^{\infty }\iiint
d^{3}k\left[ \left\{ P_{\Vert }^{\ast }E_{\mathbf{k},\Vert }^{\ast }J_{n}+%
\left[ P_{\bot }^{\ast }-\frac{1}{p_{\bot }}\left( \frac{k_{\Vert }v_{\Vert }%
}{\omega ^{\ast }}-1\right) \right] E_{\mathbf{k},\bot }^{\ast }\right\}
\right.  \label{qloporig} \\
&&\left. \cdot \frac{i}{\left[ n\Omega +k_{\Vert }v_{\Vert }-\omega \right] }%
\left\{ E_{\mathbf{k},\Vert }J_{n}P_{\Vert }+E_{\mathbf{k},\bot }P_{\bot
}\right\} f\right]  \notag
\end{eqnarray}
with 
\begin{eqnarray}
P_{\Vert } &=&\frac{\partial }{\partial p_{\Vert }}-\frac{n\Omega }{\omega
v_{\bot }}\left( v_{\bot }\frac{\partial }{\partial p_{\Vert }}-v_{\Vert }%
\frac{\partial }{\partial p_{\bot }}\right)  \label{ppar} \\
P_{\bot } &=&\frac{\partial }{\partial p_{\bot }}+\frac{k_{\Vert }}{\omega }%
\left( v_{\bot }\frac{\partial }{\partial p_{\Vert }}-v_{\Vert }\frac{%
\partial }{\partial p_{\bot }}\right)  \label{pperp} \\
E_{\mathbf{k},\bot } &=&\frac{1}{\sqrt{2}}\left( E_{\mathbf{k},-}e^{i\alpha
}J_{n+1}+E_{\mathbf{k},+}e^{-i\alpha }J_{n-1}\right)  \label{eperp}
\end{eqnarray}

The electric field is decomposed into its Fourier components 
\begin{equation}
\mathbf{E}(\mathbf{x},t)=\iiint \frac{d^{3}k}{(2\pi )^{3}}e^{i\mathbf{k}%
\cdot \mathbf{x}-i\omega t}\mathbf{E}_{\mathbf{k}}\left( \mathbf{k}\right)
\end{equation}
which are projected on the rotating field frame 
\begin{eqnarray}
E_{\mathbf{k},\pm } &=&\frac{1}{\sqrt{2}}\left( E_{\mathbf{k},x}\pm iE_{%
\mathbf{k},y}\right)  \label{ekpm} \\
E_{\mathbf{k},\Vert } &=&E_{\mathbf{k},z}
\end{eqnarray}

The wave vector is expressed in cylindrical coordinates as 
\begin{eqnarray}
k_{x} &=&k_{\bot }\cos \alpha \\
k_{y} &=&k_{\bot }\sin \alpha  \notag \\
k_{z} &=&k_{\Vert }  \notag
\end{eqnarray}
and the argument of the Bessel functions is $k_{\bot }v_{\bot }/\Omega $.
The relativistic cyclotron frequency is 
\begin{equation}
\Omega =\frac{q_{e}B}{\gamma m_{e}}
\end{equation}

\section{Expression as the Divergence of a Flux}

\subsection{In Spherical Coordinates}

After some manipulations (Appendix \ref{fluxexp}), the quasilinear operator (%
\ref{qloporig}) can be expressed in a divergence form in the $\left( p,\xi
\right) $ momentum space 
\begin{equation}
Q(f)=-\nabla \cdot \mathbf{S}^{\text{RF}}=-\frac{1}{p^{2}}\frac{\partial }{%
\partial p}\left( p^{2}S_{p}^{\text{RF}}\right) +\frac{1}{p}\frac{\partial }{%
\partial \xi }\left( \sqrt{1-\xi ^{2}}S_{\xi }^{\text{RF}}\right)
\label{qlop}
\end{equation}
where the flux has a purely diffusive form 
\begin{equation}
\mathbf{S}^{\text{RF}}=-\mathbb{D}^{\text{RF}}\cdot \nabla _{\mathbf{p}}f
\end{equation}

The components of the diffusion tensor are given by 
\begin{equation}
\mathbb{D}^{\text{RF}}=\left( 
\begin{array}{cc}
D_{pp}^{\text{RF}} & D_{p\xi }^{\text{RF}} \\ 
D_{\xi p}^{\text{RF}} & D_{\xi \xi }^{\text{RF}}
\end{array}
\right)
\end{equation}
with 
\begin{eqnarray}
D_{pp}^{\text{RF}} &=&\sum_{n=-\infty }^{+\infty }\sum_{b}(1-\xi
^{2})D_{b,n}^{\text{RF}}(p,\xi )  \label{dtensor} \\
D_{p\xi }^{\text{RF}} &=&\sum_{n=-\infty }^{+\infty }\sum_{b}-\frac{\sqrt{%
1-\xi ^{2}}}{\xi }\left[ 1-\xi ^{2}-\frac{n\Omega }{\omega _{b}}\right]
D_{b,n}^{\text{RF}}(p,\xi )  \notag \\
D_{\xi p}^{\text{RF}} &=&\sum_{n=-\infty }^{+\infty }\sum_{b}-\frac{\sqrt{%
1-\xi ^{2}}}{\xi }\left[ 1-\xi ^{2}-\frac{n\Omega }{\omega _{b}}\right]
D_{b,n}^{\text{RF}}(p,\xi )  \notag \\
D_{\xi \xi }^{\text{RF}} &=&\sum_{n=-\infty }^{+\infty }\sum_{b}\frac{1}{\xi
^{2}}\left[ 1-\xi ^{2}-\frac{n\Omega }{\omega _{b}}\right] ^{2}D_{b,n}^{%
\text{RF}}(p,\xi )  \notag
\end{eqnarray}
where the contributions from the rays are formally separated. Note that the
rays may describe several RF beams, so that their frequency $\omega _{b}$
can be different.

The gradient in momentum space is the vector 
\begin{equation}
\nabla _{\mathbf{p}}=\left( 
\begin{array}{l}
\nabla _{p} \\ 
\nabla _{\xi }
\end{array}
\right)
\end{equation}
with 
\begin{eqnarray}
\nabla _{p} &=&\frac{\partial f}{\partial p} \\
\nabla _{\xi } &=&-\frac{\sqrt{1-\xi ^{2}}}{p}\frac{\partial f}{\partial \xi 
}
\end{eqnarray}

The quasilinear diffusion coefficient in (\ref{dtensor}), describing the
interaction of the electrons with a given beam $b$ at an harmonic $n$, is
given by 
\begin{equation}
D_{b,n}^{\text{RF}}(p,\xi )=\lim_{V\rightarrow \infty }\frac{\pi }{V}\iiint 
\frac{d^{3}k}{(2\pi )^{3}}D_{\mathbf{k}}^{b,\left( n\right) }\delta (\omega
_{b}-k_{\Vert }v_{\Vert }-n\Omega )  \label{drf}
\end{equation}
where $D_{\mathbf{k}}^{b,(n)}$ accounts for the polarization and the
intensity of the RF wave. It is given by 
\begin{equation}
D_{\mathbf{k}}^{b,(n)}=e^{2}\left| \frac{1}{\sqrt{2}}E_{\mathbf{k}%
,b,+}e^{-i\alpha }J_{n-1}\left( \frac{k_{\bot ,b}v_{\bot }}{\Omega }\right) +%
\frac{1}{\sqrt{2}}E_{\mathbf{k},b,-}e^{+i\alpha }J_{n+1}\left( \frac{k_{\bot
,b}v_{\bot }}{\Omega }\right) +\frac{p_{\Vert }}{p_{\bot }}E_{\mathbf{k}%
,b,\Vert }J_{n}\left( \frac{k_{\bot ,b}v_{\bot }}{\Omega }\right) \right|
^{2}  \label{dk2}
\end{equation}

\subsection{In Cylindrical Coordinates}

The FP equation is usually solved numerically in spherical coordinates,
because of the spherical structure of the collisional operator. However, in
many case the RF QL Operator has a more cylindrical structure, which is
useful in order to get a physical picture of the wave-particle interaction.
The expression of the divergence in cylindrical coordinates is 
\begin{equation}
\nabla \cdot \mathbf{S}=\frac{1}{p_{\bot }^{2}}\frac{\partial }{\partial
p_{\bot }}\left( p_{\bot }S_{\bot }\right) +\frac{\partial }{\partial
p_{\Vert }}\left( S_{\Vert }\right)  \label{divcyl}
\end{equation}
If the fluxes have a diffusive form 
\begin{equation}
\mathbf{S}=-\mathbb{D}\cdot \nabla f=-\left( 
\begin{array}{cc}
D_{\bot \bot } & D_{\bot \Vert } \\ 
D_{\Vert \bot } & D_{\Vert \Vert }
\end{array}
\right) \cdot \left( 
\begin{array}{c}
\partial f/\partial p_{\bot } \\ 
\partial f/\partial p_{\Vert }
\end{array}
\right)
\end{equation}
the cylindrical tensor elements are related to the spherical tensor elements
by 
\begin{equation}
\left( 
\begin{array}{c}
D_{\bot \bot } \\ 
D_{\bot \Vert } \\ 
D_{\Vert \bot } \\ 
D_{\Vert \Vert }
\end{array}
\right) =\left( 
\begin{array}{cccc}
\left( 1-\xi ^{2}\right) & \xi \sqrt{1-\xi ^{2}} & \xi \sqrt{1-\xi ^{2}} & 
\xi ^{2} \\ 
\xi \sqrt{1-\xi ^{2}} & -\left( 1-\xi ^{2}\right) & \xi ^{2} & -\xi \sqrt{%
1-\xi ^{2}} \\ 
\xi \sqrt{1-\xi ^{2}} & \xi ^{2} & -\left( 1-\xi ^{2}\right) & -\xi \sqrt{%
1-\xi ^{2}} \\ 
\xi ^{2} & -\xi \sqrt{1-\xi ^{2}} & -\xi \sqrt{1-\xi ^{2}} & \left( 1-\xi
^{2}\right)
\end{array}
\right) \cdot \left( 
\begin{array}{c}
D_{pp} \\ 
D_{p\xi } \\ 
D_{\xi p} \\ 
D_{\xi \xi }
\end{array}
\right)
\end{equation}

Applying this transformation to the RF diffusion tensor elements above gives 
\begin{eqnarray}
D_{\bot \bot }^{\text{RF}} &=&\sum_{n=-\infty }^{+\infty }\sum_{b}\left[ 
\frac{n\Omega }{\omega _{b}}\right] ^{2}D_{b,n}^{\text{RF}}(p,\xi )
\label{dcyl} \\
D_{\bot \Vert }^{\text{RF}} &=&\sum_{n=-\infty }^{+\infty }\sum_{b}\frac{%
\sqrt{1-\xi ^{2}}}{\xi }\frac{n\Omega }{\omega _{b}}\left[ 1-\frac{n\Omega }{%
\omega _{b}}\right] D_{b,n}^{\text{RF}}(p,\xi )  \notag \\
D_{\Vert \bot }^{\text{RF}} &=&\sum_{n=-\infty }^{+\infty }\sum_{b}\frac{%
\sqrt{1-\xi ^{2}}}{\xi }\frac{n\Omega }{\omega _{b}}\left[ 1-\frac{n\Omega }{%
\omega _{b}}\right] D_{b,n}^{\text{RF}}(p,\xi )  \notag \\
D_{\Vert \Vert }^{\text{RF}} &=&\sum_{n=-\infty }^{+\infty }\sum_{b}\frac{%
(1-\xi ^{2})}{\xi ^{2}}\left[ 1-\frac{n\Omega }{\omega _{b}}\right]
^{2}D_{b,n}^{\text{RF}}(p,\xi )  \notag
\end{eqnarray}

We see immediately that for $n=0$ (for example, LHCD) the QL diffusion is
strictly in the parallel direction. Also, at a cyclotron harmonic, where $%
\omega _{b}=n\Omega $, the diffusion is only perpendicular (not considering
relativistic effects).

\section{Bounce Averaging of the RF Quasilinear Operator}

The RF Quasilinear Operator is bounce averaged in order to revove the
dependence on the poloidal angle. The bounce averaging operation in both
toroidal and poloidal directions is defined as 
\begin{equation}
\left\{ \mathcal{A}\right\} =\frac{1}{\lambda \left( \xi _{0}\right) }\left[ 
\frac{1}{2}\sum_{\sigma =\pm 1}\right] _{T}\int_{-\theta _{c}}^{\theta _{c}}%
\frac{d\theta }{2\pi }\int_{\varsigma _{0}-\varsigma _{c}}^{\varsigma
_{0}+\varsigma _{c}}\frac{d\varsigma }{2\pi }\frac{\xi _{0}}{\xi }\mathcal{A}
\end{equation}
where $\xi _{0}$ is the value of $\xi $ and $\varsigma _{0}$ is the toroidal
position of the particle as it crosses the low field side midplane. More
generally, the subscript $_{0}$ refers to quantities evaluated on the low
field side midplane. The sign of $\xi $ and $\xi _{0}$ is

\begin{equation}
\sigma =\frac{\xi _{0}}{\left| \xi _{0}\right| }
\end{equation}
and the sum over $\sigma $\ applies only to trapped particles.

The poloidal integration boundary is 
\begin{eqnarray}
\theta _{c} &=&\pi \text{ for passing particles} \\
\theta _{c} &=&\theta _{T}\text{ for trapped particles}
\end{eqnarray}
where 
\begin{equation}
\theta _{T}=\cos ^{-1}\left( 1-\frac{2\xi _{0}^{2}}{\xi _{0T}^{2}}\right)
\end{equation}
is the poloidal angle at the turning point. The trapping occurs for
electrons with $\left| \xi _{0}\right| <\xi _{0T}\ $with 
\begin{equation}
\xi _{0T}^{2}=\frac{2\epsilon }{1+\epsilon }
\end{equation}
where 
\begin{equation}
\epsilon =\frac{r}{R_{0}}
\end{equation}
is the local inverse aspeect ratio. The normalized bounce time is 
\begin{equation}
\lambda \left( \xi _{0}\right) =\int_{-\theta _{c}}^{\theta _{c}}\frac{%
d\theta }{2\pi }\frac{\xi _{0}}{\xi }
\end{equation}

We have

\begin{equation}
\frac{d\varsigma }{d\theta }=\frac{B_{\varsigma }\left( r,\theta \right) }{%
B_{\theta }\left( r,\theta \right) }=\frac{R_{0}}{r}q\left( r\right) =\frac{%
q\left( r\right) }{\epsilon }
\end{equation}
so that 
\begin{equation}
\varsigma _{c}=\left\{ 
\begin{array}{cc}
\pi & \text{ if }\theta _{c}=\pi \text{ or }\theta _{T}/\pi >\epsilon /q \\ 
\theta _{T}q/\epsilon & \text{ if }\theta _{c}<\pi \text{ and }\theta
_{T}/\pi <\epsilon /q
\end{array}
\right.
\end{equation}

In the second case, particles are so trapped that the extent of their orbit
in the toroidal direction is less that $2\pi $.

Defining the ratio of magnetic fields 
\begin{equation}
\Psi \left( r,\theta \right) =\frac{B\left( r,\theta \right) }{B_{0}\left(
r,\theta \right) }=\frac{R_{0}}{R}=\frac{1+\epsilon }{1+\epsilon \cos \theta 
}  \label{psidef}
\end{equation}
we have 
\begin{equation}
\xi =\sigma \sqrt{1-\Psi \left( r,\theta \right) \left( 1-\xi
_{0}^{2}\right) }
\end{equation}

The bounce-averaged RF quasilinear operator conserves its divergence form in
the in the $\left( p,\xi _{0}\right) $ momentum space 
\begin{equation}
\left\{ Q(f_{0})\right\} =\left\{ -\nabla \cdot \mathbf{S}^{\text{RF}%
}\right\} =-\frac{1}{p^{2}}\frac{\partial }{\partial p}\left( p^{2}S_{p}^{%
\text{RF}(0)}\right) +\frac{1}{\lambda \left( \xi _{0}\right) p}\frac{%
\partial }{\partial \xi _{0}}\left( \sqrt{1-\xi _{0}^{2}}\lambda \left( \xi
_{0}\right) S_{\xi _{0}}^{\text{RF}(0)}\right)
\end{equation}
where the flux in momentum space is still purely diffusive 
\begin{equation}
\mathbf{S}^{\text{RF}(0)}=-\mathbb{D}^{\text{RF}(0)}\cdot \nabla _{\mathbf{p}%
}^{(0)}f^{(0)}
\end{equation}
where the component of the diffusion tensor are given by 
\begin{eqnarray}
D_{pp}^{\text{RF}(0)} &=&\sum_{n=-\infty }^{+\infty }\sum_{b}(1-\xi
_{0}^{2})D_{b,n}^{\text{RF}(0)}(p,\xi _{0}) \\
D_{p\xi }^{\text{RF}(0)} &=&\sum_{n=-\infty }^{+\infty }\sum_{b}-\frac{\sqrt{%
1-\xi _{0}^{2}}}{\xi _{0}}\left[ 1-\xi _{0}^{2}-\frac{n\Omega _{0}}{\omega
_{b}}\right] D_{b,n}^{\text{RF}(0)}(p,\xi _{0})  \notag \\
D_{\xi p}^{\text{RF}(0)} &=&\sum_{n=-\infty }^{+\infty }\sum_{b}-\frac{\sqrt{%
1-\xi _{0}^{2}}}{\xi _{0}}\left[ 1-\xi _{0}^{2}-\frac{n\Omega _{0}}{\omega
_{b}}\right] D_{b,n}^{\text{RF}(0)}(p,\xi _{0})  \notag \\
D_{\xi \xi }^{\text{RF}(0)} &=&\sum_{n=-\infty }^{+\infty }\sum_{b}\frac{1}{%
\xi _{0}^{2}}\left[ 1-\xi _{0}^{2}-\frac{n\Omega _{0}}{\omega _{b}}\right]
^{2}D_{b,n}^{\text{RF}(0)}(p,\xi _{0})  \notag
\end{eqnarray}
with the cyclotron frequency taken at the LFS midplane 
\begin{equation}
\Omega _{0}=\frac{q_{e}B_{0}}{\gamma m_{e}}=\frac{\Omega }{\Psi }
\label{omega0}
\end{equation}

The gradient in momentum space is now the vector 
\begin{equation}
\nabla _{\mathbf{p}}^{(0)}=\left( 
\begin{array}{l}
\nabla _{p} \\ 
\nabla _{\xi _{0}}
\end{array}
\right)
\end{equation}
with 
\begin{eqnarray}
\nabla _{p} &=&\frac{\partial f}{\partial p} \\
\nabla _{\xi } &=&-\frac{\sqrt{1-\xi _{0}^{2}}}{p}\frac{\partial f}{\partial
\xi _{0}}
\end{eqnarray}

This formulation for the bounce-averaged RF quasilinear operator conserves
the structure of the operator (\ref{qlop}), and a bounce-averaged
quasilinear diffusion coefficient is defined: 
\begin{equation}
D_{b,n}^{\text{RF}(0)}(p,\xi _{0})=\left\{ \Psi D_{b,n}^{\text{RF}}(p,\xi
)\right\}  \label{dbounce}
\end{equation}

\section{Diffusion coefficient for a Plane Wave\label{beam}}

\subsection{Integration in $k$-space}

The quasilinear diffusion coefficient (\ref{drf}), describing the
interaction of the electrons with a given beam $b$ at an harmonic $n$, is
rewritten as 
\begin{equation}
D_{b,n}^{\text{RF}}(p,\xi )=\lim_{V\rightarrow \infty }\frac{q_{e}^{2}\pi }{V%
}\iiint \frac{d^{3}k}{(2\pi )^{3}}\left| \mathbf{E}_{b\mathbf{k}}\right|
^{2}\left| \Theta _{\mathbf{k}}^{b,n}\right| ^{2}\delta \left( \omega
_{b}-k_{\Vert }v_{\Vert }-n\Omega \right)  \label{dn}
\end{equation}
where we define 
\begin{equation}
\Theta _{\mathbf{k}}^{b,n}=\frac{1}{\sqrt{2}}e_{b\mathbf{k},+}e^{-i\alpha
}J_{n-1}\left( \frac{k_{\bot }v_{\bot }}{\Omega }\right) +\frac{1}{\sqrt{2}}%
e_{b\mathbf{k},-}e^{+i\alpha }J_{n+1}\left( \frac{k_{\bot }v_{\bot }}{\Omega 
}\right) +\frac{p_{\Vert }}{p_{\bot }}e_{b\mathbf{k},\Vert }J_{n}\left( 
\frac{k_{\bot }v_{\bot }}{\Omega }\right)  \label{thetak}
\end{equation}
and the polarization vector 
\begin{equation}
\mathbf{e}_{b\mathbf{k}}=\frac{\mathbf{E}_{b\mathbf{k}}}{\left| \mathbf{E}_{b%
\mathbf{k}}\right| }  \label{polvec}
\end{equation}
with the components 
\begin{eqnarray}
e_{b\mathbf{k},+} &=&\frac{E_{b\mathbf{k},+}}{\left| \mathbf{E}_{b\mathbf{k}%
}\right| }=\frac{E_{b\mathbf{k},x}+iE_{b\mathbf{k},y}}{\sqrt{2}\left| 
\mathbf{E}_{b\mathbf{k}}\right| } \\
e_{b\mathbf{k},-} &=&\frac{E_{b\mathbf{k},-}}{\left| \mathbf{E}_{b\mathbf{k}%
}\right| }=\frac{E_{b\mathbf{k},x}-iE_{b\mathbf{k},y}}{\sqrt{2}\left| 
\mathbf{E}_{b\mathbf{k}}\right| } \\
e_{b\mathbf{k},\Vert } &=&\frac{E_{b\mathbf{k},\Vert }}{\left| \mathbf{E}_{b%
\mathbf{k}}\right| }=\frac{E_{b\mathbf{k},z}}{\left| \mathbf{E}_{b\mathbf{k}%
}\right| }
\end{eqnarray}

We electric field associated with a given ray is described by a plane wave,
with given frequency $\omega _{b}$ and wave number $\mathbf{k}_{b}$: 
\begin{equation}
\mathbf{E}_{b}\left( \mathbf{x},t\right) =\text{Re}\left[ \mathbf{E}%
_{b0}e^{i\left( \mathbf{k}_{b}\cdot \mathbf{x}-\omega _{b}t\right) }\right]
\end{equation}

We then get (See Appendix \ref{plane}) 
\begin{equation}
D_{b,n}^{\text{RF}}(p,\xi )=q_{e}^{2}\pi \frac{\left| \mathbf{E}_{b0}\right|
^{2}}{2}\left| \Theta _{\mathbf{k}}^{b,n}\left( \mathbf{k}_{b}\right)
\right| ^{2}\delta \left( \omega _{b}-k_{\Vert b}v_{\Vert }-n\Omega \right)
\end{equation}
which can be rewritten as\qquad \qquad \qquad \qquad 
\begin{eqnarray}
D_{b,n}^{\text{RF}}(p,\xi ) &=&q_{e}^{2}\pi \frac{\left| \mathbf{E}%
_{b0}\right| ^{2}}{2}\left| \Theta _{\mathbf{k}}^{b,n}\left( \mathbf{k}%
_{b}\right) \right| ^{2}\frac{1}{\left| v_{\Vert }\right| }\delta \left(
k_{\Vert b}-\frac{\omega _{b}}{v_{\Vert }}+\frac{n\Omega }{v_{\Vert }}\right)
\notag \\
&=&q_{e}^{2}\pi \frac{\left| \mathbf{E}_{b0}\right| ^{2}}{2}\left| \Theta _{%
\mathbf{k}}^{b,n}\left( \mathbf{k}_{b}\right) \right| ^{2}\frac{1}{\left|
v_{\Vert }\right| }\frac{c}{\omega _{b}}\delta \left( N_{\Vert b}-N_{\Vert 
\text{res}}\right)  \label{diffcoef}
\end{eqnarray}
with 
\begin{eqnarray}
N_{\Vert b} &=&\frac{k_{\Vert b}c}{\omega _{b}} \\
N_{\Vert \text{res}} &=&\frac{c}{v_{\Vert }}\left( 1-\frac{n\Omega }{\omega
_{b}}\right)  \notag \\
&=&\frac{1}{\beta _{Te}}\frac{p_{Te}}{p_{\Vert }}\left( \gamma -\frac{%
n^{\prime }\omega _{ce}}{\omega _{b}}\right)
\end{eqnarray}
where 
\begin{eqnarray}
\Omega &=&\frac{q_{e}B}{\gamma m_{e}}=\frac{-\omega _{ce}}{\gamma }\text{
with }\omega _{ce}=\frac{eB}{m_{e}} \\
n^{\prime } &=&-n \\
\beta _{Te} &=&\frac{p_{Te}}{m_{e}c}=\frac{v_{Te}}{c}=\sqrt{\frac{kT_{e}}{%
m_{e}c^{2}}}
\end{eqnarray}

\subsection{Incident Energy Flow Density and Relation to the Electric Field}

The time-averaged energy flow density in the beam is in general given by 
\cite{sti92} 
\begin{equation}
\overline{\mathbf{s}_{b}}=\mathbf{P}_{b}+\mathbf{T}_{b}  \label{sb}
\end{equation}
where 
\begin{equation}
\mathbf{P}_{b}=\frac{1}{2}\text{Re}\left[ \widetilde{\mathbf{E}}_{b}\times 
\widetilde{\mathbf{H}}_{b}^{\ast }\right]
\end{equation}
is the flow of electromagnetic energy, or Poynting vector, and 
\begin{equation}
\mathbf{T}_{b}=-\frac{\epsilon _{0}\omega _{b}}{4}\widetilde{\mathbf{E}}%
_{b}^{\ast }\cdot \frac{\partial \mathbb{K}^{H}}{\partial \mathbf{k}}\cdot 
\widetilde{\mathbf{E}}_{b}
\end{equation}
is the kinetic energy flow density. $\widetilde{\mathbf{E}}_{b}$ is the
complex form of the electric field 
\begin{equation}
\widetilde{\mathbf{E}}_{b}=\mathbf{E}_{b0}e^{i\mathbf{k\cdot x}-i\omega t}
\end{equation}
and $\mathbb{K}^{H}$ is the Hermitian part of the dielectric tensor.

The energy flow density (\ref{sb}) can be formally rewritten as 
\begin{equation}
\overline{\mathbf{s}_{b}}=\frac{\epsilon _{0}c}{2}\left| \mathbf{E}%
_{b0}\right| ^{2}\mathbf{\Phi }_{b}
\end{equation}
where $\mathbf{\Phi }_{b}$ is a non dimensional vector defined as 
\begin{equation}
\mathbf{\Phi }_{b}=\mathbf{\Phi }_{bP}+\mathbf{\Phi }_{bT}  \label{phib2}
\end{equation}
with 
\begin{equation}
\mathbf{\Phi }_{bP}=\frac{\mu _{0}c}{\left| \mathbf{E}_{b0}\right| ^{2}}%
\text{Re}\left[ \widetilde{\mathbf{E}}_{b}\times \widetilde{\mathbf{H}}%
_{b}^{\ast }\right]  \label{phib}
\end{equation}
and 
\begin{equation}
\mathbf{\Phi }_{bT}=\mathbf{-}\frac{1}{2}\mathbf{e}_{b}^{\ast }\cdot \frac{%
\partial \mathbb{K}^{H}}{\partial \mathbf{N}}\cdot \mathbf{e}_{b}
\label{phit}
\end{equation}

In vacuum, $\mathbf{\Phi }_{b}=\widehat{k}$ , unit vector in the direction
of propagation, and $\mathbf{\Phi }_{bT}=0$. The vector $\mathbf{e}_{b}$ is
the polarization vector, defined in (\ref{polvec}).

The incident power flux on a flux-surface is therefore given by 
\begin{equation}
\overline{\mathbf{s}_{b}}_{\text{inc}}=\frac{\epsilon _{0}c}{2}\left| 
\mathbf{E}_{b0}\right| ^{2}\left| \mathbf{\Phi }_{b}\cdot \widehat{n}\right|
\label{sbinc}
\end{equation}
where $\widehat{n}$ is the\ local vector normal to the flux-surface.

\subsection{Approximation of a Narrow Beam}

We consider a ray which crosses the flux surface at the respective poloidal
and toroidal locations $\theta _{b}$ and $\varsigma _{b}$.\ If we make the
asumption that the extension of the beam power on the flux-surface is very
small in both the poloidal and the toroidal directions, we can approximate 
\begin{equation}
\overline{\mathbf{s}_{b}}_{\text{inc}}\simeq \frac{4\pi ^{2}}{A_{b}}\delta
\left( \theta -\theta _{b}\right) \delta \left( \varsigma -\varsigma
_{b}\right) P_{b,\text{inc}}
\end{equation}

Note that this asumption is in contradication with the plane-wave
decomposition of the RF beam, since plane waves are essentially present in
the entire configuration space. However, these plane waves interfere
destructively everywhere except in the beam region, which can be limited to
a narrow poloidal and toroidal extention.

$A_{b}$ is to be determined such that $P_{\text{inc}}$ is the total incident
power on the flux-surface 
\begin{equation}
\iint dS\;\overline{\mathbf{s}_{b}}_{\text{inc}}=P_{b,\text{inc}}
\end{equation}
which gives 
\begin{eqnarray}
1 &=&\int_{0}^{2\pi }rd\theta \;\int_{0}^{2\pi }Rd\varsigma \;\frac{4\pi ^{2}%
}{A_{b}}\delta \left( \theta -\theta _{b}\right) \delta \left( \varsigma
-\varsigma _{b}\right)  \notag \\
&=&\frac{4\pi ^{2}rR_{c}}{A_{b}}\int_{0}^{2\pi }d\theta \;\frac{\Psi _{c}}{%
\Psi }\delta \left( \theta -\theta _{b}\right)  \notag \\
&=&\frac{4\pi ^{2}rR_{c}}{A_{b}}\;\frac{\Psi _{c}}{\Psi _{b}}
\end{eqnarray}
where we used (\ref{psidef}) so that 
\begin{equation}
R\left( r,\theta \right) =R_{c}\frac{\Psi _{c}}{\Psi \left( r,\theta \right) 
}  \label{psi}
\end{equation}
and defined 
\begin{eqnarray}
\Psi _{b} &\equiv &\Psi \left( r,\theta _{b}\right) \\
\Psi _{c} &\equiv &\Psi \left( r,\pi /2\right) =\frac{B_{c}}{B_{0}}
\end{eqnarray}
where $B_{c}$ is the field on axis ($R=R_{c}$). For concentric cylindrical
flux-surfaces, we have 
\begin{eqnarray}
\Psi _{b} &=&\frac{1+\epsilon }{1+\epsilon \cos \theta _{b}} \\
\Psi _{c} &=&1+\epsilon
\end{eqnarray}

Finally 
\begin{equation}
A_{b}=4\pi ^{2}rR_{c}\frac{\Psi _{c}}{\Psi _{b}}  \label{psi0}
\end{equation}
and the incident power flux is 
\begin{equation}
\overline{\mathbf{s}_{b}}_{\text{inc}}\simeq \frac{1}{rR_{c}}\frac{\Psi _{b}%
}{\Psi _{c}}\delta \left( \theta -\theta _{b}\right) \delta \left( \varsigma
-\varsigma _{b}\right) P_{b,\text{inc}}  \label{sbinc2}
\end{equation}

Using (\ref{sbinc}) and (\ref{sbinc2}), we find 
\begin{equation}
\frac{\left| \mathbf{E}_{b0}\right| ^{2}}{2}=\frac{1}{\epsilon _{0}c}\frac{1%
}{\left| \mathbf{\Phi }_{b}\cdot \widehat{n}\right| }\frac{1}{rR_{c}}\frac{%
\Psi _{b}}{\Psi _{c}}\delta \left( \theta -\theta _{b}\right) \delta \left(
\varsigma -\varsigma _{b}\right) P_{b,\text{inc}}  \label{e0}
\end{equation}

\subsection{Normalized Diffusion Coefficient}

The diffusion coefficient is usually normalized to the collisional diffusion
coefficient $\nu _{e}p_{Te}^{2}$, thus defining 
\begin{equation}
\overline{D}_{b,n}^{\text{RF}}(p,\xi )\equiv \frac{D_{b,n}^{\text{RF}}(p,\xi
)}{\nu _{e}p_{Te}^{2}}  \label{difcoeffnorm}
\end{equation}
with 
\begin{equation}
\nu _{e}=\frac{e^{4}n_{e}\ln \Lambda }{4\pi \epsilon
_{0}^{2}m_{e}^{2}v_{Te}^{3}}
\end{equation}

Using (\ref{diffcoef}) and (\ref{e0}), the diffusion coefficient (\ref
{difcoeffnorm}) becomes 
\begin{equation}
\overline{D}_{b,n}^{\text{RF}}(p,\xi )=2\pi \delta \left( \theta -\theta
_{b}\right) 2\pi \delta \left( \varsigma -\varsigma _{b}\right) \left|
\Theta _{\mathbf{k}}^{b,n}\left( \mathbf{k}_{b}\right) \right| ^{2}\frac{%
\gamma p_{Te}}{p\left| \xi \right| }\overline{D}_{b,n,0}^{\text{RF}}\delta
\left( N_{\Vert b}-N_{\Vert \text{res}}\right)  \label{difcoeffnorm2}
\end{equation}
with 
\begin{eqnarray}
\overline{D}_{b,n,0}^{\text{RF}} &=&\frac{1}{\nu _{e}p_{Te}^{2}}\frac{%
q_{e}^{2}\pi }{\epsilon _{0}c}\frac{P_{b,\text{inc}}}{4\pi ^{2}rR_{c}}\frac{%
\Psi _{b}}{\Psi _{c}}\frac{1}{\left| \mathbf{\Phi }_{b}\cdot \widehat{n}%
\right| }\frac{1}{v_{Te}}\frac{c}{\omega _{b}} \\
&=&\frac{4\pi }{m_{e}\omega _{pe}^{2}\ln \Lambda }\frac{\pi }{\omega _{b}}%
\frac{\Psi _{b}}{\Psi _{c}}\frac{1}{\left| \mathbf{\Phi }_{b}\cdot \widehat{n%
}\right| }\frac{P_{b,\text{inc}}}{4\pi ^{2}rR_{c}}
\end{eqnarray}
where we introduce the plasma frequency 
\begin{equation}
\omega _{pe}=\sqrt{\frac{e^{2}n_{e}}{\epsilon _{0}m_{e}}}
\end{equation}

\subsection{Bounce-Averaging of the Diffusion Coefficient}

The bounce-averaged diffusion coefficient is obtained from (\ref
{difcoeffnorm2}) and (\ref{dbounce}) as 
\begin{equation}
\overline{D}_{b,n}^{\text{RF}(0)}(p,\xi _{0})=\left\{ \Psi 2\pi \delta
\left( \theta -\theta _{b}\right) 2\pi \delta \left( \varsigma -\varsigma
_{b}\right) \left| \Theta _{\mathbf{k}}^{b,n}\left( \mathbf{k}_{b}\right)
\right| ^{2}\frac{\gamma p_{Te}}{p\left| \xi \right| }\overline{D}_{b,n,0}^{%
\text{RF}}\delta \left( N_{\Vert b}-N_{\Vert \text{res}}\right) \right\}
\end{equation}

Recalling the expression of bounce-averaging 
\begin{equation}
\left\{ \mathcal{A}\right\} =\frac{1}{\lambda \left( \xi _{0}\right) }\left[ 
\frac{1}{2}\sum_{\sigma =\pm 1}\right] _{T}\int_{-\theta _{c}}^{\theta _{c}}%
\frac{d\theta }{2\pi }\int_{\varsigma _{0}-\varsigma _{c}}^{\varsigma
_{0}+\varsigma _{c}}\frac{d\varsigma }{2\pi }\frac{\xi _{0}}{\xi }\mathcal{A}
\end{equation}
we do the integrals, which gives 
\begin{eqnarray}
\overline{D}_{b,n}^{\text{RF}(0)}(p,\xi _{0}) &=&\frac{\gamma p_{Te}}{%
p\left| \xi _{0}\right| }\frac{\Psi _{b}^{2}}{\Psi _{c}}\frac{\xi _{0}^{2}}{%
\xi _{b}^{2}}\left| \Theta _{\mathbf{k},b}^{b,n}\left( \mathbf{k}_{b}\right)
\right| ^{2}\overline{D}_{b,n,0}^{\text{RF},b}\frac{1}{\lambda \left( \xi
_{0}\right) } \\
&&H\left( \varsigma _{c}-\left| \varsigma _{b}-\varsigma _{0}\right| \right)
H\left( \theta _{c}-\left| \theta _{b}\right| \right) \left[ \frac{1}{2}%
\sum_{\sigma =\pm 1}\right] _{T}\delta \left( N_{\Vert b}-N_{\Vert \text{res}%
,b}\right)
\end{eqnarray}
where we define: 
\begin{eqnarray}
\overline{D}_{b,n,0}^{\text{RF},b} &=&\frac{1}{\nu _{e}p_{Te}^{2}}\frac{%
q_{e}^{2}\pi }{\epsilon _{0}\omega _{b}}\frac{P_{b,\text{inc}}}{4\pi
^{2}rR_{c}}\frac{\Psi _{b}}{\Psi _{c}}\frac{1}{\left| \mathbf{\Phi }%
_{b}\cdot \widehat{n}_{b}\right| }\frac{1}{v_{Te}} \\
&=&\frac{4\pi }{m_{e}\omega _{pe}^{2}\ln \Lambda }\frac{\pi }{\omega _{b}}%
\frac{\Psi _{b}}{\Psi _{c}}\frac{1}{\left| \mathbf{\Phi }_{b}\cdot \widehat{n%
}_{b}\right| }\frac{P_{b,\text{inc}}}{4\pi ^{2}rR_{c}}
\end{eqnarray}
\begin{eqnarray}
\widehat{n}_{b} &=&\widehat{n}\left( \theta _{b}\right) \\
N_{\Vert \text{res},b} &=&\frac{1}{\beta _{Te}}\frac{p_{Te}}{p\xi _{b}}%
\left( \gamma -\frac{n^{\prime }\Psi _{b}\omega _{ce,0}}{\omega _{b}}\right)
\\
\xi _{b} &=&\sigma \sqrt{1-\Psi _{b}\left( 1-\xi _{0}^{2}\right) }
\end{eqnarray}
and 
\begin{equation}
\Theta _{\mathbf{k},b}^{b,n}\left( \mathbf{k}_{b}\right) =\frac{1}{\sqrt{2}}%
\varepsilon _{b\mathbf{k},+}e^{-i\alpha }J_{n-1}\left( z_{b}\right) +\frac{1%
}{\sqrt{2}}\varepsilon _{b\mathbf{k},-}e^{+i\alpha }J_{n+1}\left(
z_{b}\right) +\frac{\xi _{b}}{\sqrt{\Psi _{b}\left( 1-\xi _{0}^{2}\right) }}%
\varepsilon _{b\mathbf{k},\Vert }J_{n}\left( z_{b}\right)  \label{zb}
\end{equation}
with 
\begin{equation}
z_{b}=\frac{k_{\bot b}v_{\bot b}}{\Omega _{b}}=-N_{\bot b}\frac{p}{m_{e}c}%
\sqrt{1-\xi _{b}^{2}}\frac{\omega _{b}}{\omega _{ce}}=-N_{\bot b}\frac{p}{%
m_{e}c}\frac{\sqrt{1-\xi _{0}^{2}}}{\sqrt{\Psi _{b}}}\frac{\omega _{b}}{%
\omega _{ce,0}}  \label{zb2}
\end{equation}
where we used (\ref{omega0}).

Averaging over $\varsigma _{0}$ gives 
\begin{equation}
\overline{D}_{b,n}^{\text{RF}(0)}(p,\xi _{0})=\frac{\gamma p_{Te}}{p\left|
\xi _{0}\right| }\left| \Theta _{\mathbf{k},b}^{b,n}\left( \mathbf{k}%
_{b}\right) \right| ^{2}\overline{D}_{b,n,0}^{\text{RF},b}\Psi _{b}\frac{\xi
_{0}^{2}}{\xi _{b}^{2}}\frac{H_{\text{T}}\left( \xi _{0}\right) }{\lambda
\left( \xi _{0}\right) }\left[ \frac{1}{2}\sum_{\sigma =\pm 1}\right]
_{T}\delta \left( N_{\Vert b}-N_{\Vert \text{res},b}\right)
\label{diffcoefnorm3}
\end{equation}
with 
\begin{equation}
H_{\text{T}}\left( \xi _{0}\right) \equiv \frac{\varsigma _{c}}{\pi }H\left(
\theta _{c}-\left| \theta _{b}\right| \right) =\left\{ 
\begin{array}{cc}
1 & \text{if }\left| \xi _{0}\right| >\xi _{0T}\text{ (passing electrons)}
\\ 
1 & \text{if }\left| \xi _{0}\right| <\xi _{0T}\text{ , }\theta _{T}>\left|
\theta _{b}\right| \text{ and }\theta _{T}>\epsilon \pi /q \\ 
\theta _{T}q/\epsilon & \text{if }\left| \xi _{0}\right| <\xi _{0T}\text{ , }%
\theta _{T}>\left| \theta _{b}\right| \text{ and }\theta _{T}<\epsilon \pi /q
\\ 
0 & \text{if }\left| \xi _{0}\right| <\xi _{0T}\text{ \ and }\theta
_{T}<\left| \theta _{b}\right|
\end{array}
\right.
\end{equation}

\subsection{Conclusion}

In the general case, the dispersion relation must be solved for each ray $b$
in order to determine $k_{\bot b}$ in (\ref{zb2}) as a function of $k_{\Vert
b}$ and $\omega _{b}$; then, the polarization components $\varepsilon _{b%
\mathbf{k},+}$, $\varepsilon _{b\mathbf{k},-}$ and $\varepsilon _{b\mathbf{k}%
,\Vert }$\ of the electric field must be evaluated from the elements of the
dielectric tensor. Therefore, both $\Theta _{\mathbf{k},b}^{b,n}\left( 
\mathbf{k}_{b}\right) $ (\ref{zb}) and $\mathbf{\Phi }_{b}$ (\ref{phib2})
can be determined from the dielectric tensor, knowing $k_{\Vert b}$, which
can be determined along the ray path by ray-tracing techniques.

A numerical code is required to solve the wave equation, and evaluate the
diffusion coefficient\ (\ref{diffcoefnorm3}) in the general case. However,
it is possible to obtain analytic expressions in some cases, assuming that
the waves can be treated in the cold plasma limit.

\section{Diffusion coefficient for Lower Hybrid Current Drive}

Analytic expression for characteristics (spectrum and polarization) of LH
waves can be obtained in the Cold Plasma description. These are derived in
Appendix \ref{lhw}, and the useful expressions are summarized below. Then,
an analytical expression can be obtained for the LH diffusion coefficient.

\subsection{ Cold Plasma Characteristics of LH waves (Appendix \ref{lhw})}

\subsubsection{Dispersion Relation}

The electrostatic approximation (\ref{esapprox}) is usually valid for LHW of
concern. Solving the electrostatic dispersion relation (\ref{dispreles})
gives $N_{b\perp }$ as a function of $N_{b\parallel }$ and $\omega _{b}$ (%
\ref{nperplh0}-\ref{nperplh2}): 
\begin{eqnarray}
N_{b\perp } &=&\sqrt{\frac{-K_{\parallel }}{K_{\perp }}}N_{b\parallel } 
\notag \\
&=&\frac{1}{\sqrt{\omega _{b}^{2}/\omega _{\text{LH}}^{2}-1}}\frac{\omega
_{pe}}{\omega _{pi}}N_{b\parallel }  \label{nperplh1}
\end{eqnarray}
where $\omega _{\text{LH}}$\ is the lower-hybrid frequency 
\begin{equation}
\omega _{\text{LH}}=\frac{\omega _{pi}}{\sqrt{1+\omega _{pe}^{2}/\omega
_{ce}^{2}}}
\end{equation}

We can verify that LHW encounter a resonance at the lower-hybrid frequency.
However, in most LHCD scenarios, wave frequency are sensibly higher than the
lower-hybrid frequency in order to avoid conversion to IBW. Typically, in
Alcator C-Mod, we have 
\begin{equation}
\frac{\omega }{\omega _{\text{LH}}}\sim 2\text{ or }3
\end{equation}

The cold plasma description therefore remains valid.

\subsubsection{Polarization}

The LH Waves of concern for LHCD are quasi-electrostatic, and therefore the
electric field is quasi-longitudinal ($\mathbf{E}\parallel \mathbf{N}$) and
we have\textbf{\ (}\ref{eratio}\textbf{)} 
\begin{equation}
e_{b\mathbf{k},i}=\frac{E_{b\mathbf{k},i}}{\left| \mathbf{E}_{b\mathbf{k}%
}\right| }\simeq \frac{k_{b,i}}{\left| \mathbf{k}_{b}\right| }
\end{equation}
for any component $i$, so that the polarization elements become 
\begin{eqnarray}
e_{b\mathbf{k},+} &\equiv &\frac{E_{b\mathbf{k},x}+iE_{b\mathbf{k},y}}{\sqrt{%
2}\left| \mathbf{E}_{b\mathbf{k}}\right| }=\frac{k_{b\bot }}{\sqrt{2}k_{b}}%
\left( \cos \alpha +i\sin \alpha \right) =\frac{k_{b\bot }}{\sqrt{2}k_{b}}%
e^{+i\alpha }  \label{pollh} \\
e_{b\mathbf{k},-} &\equiv &\frac{E_{b\mathbf{k},x}-iE_{b\mathbf{k},y}}{\sqrt{%
2}\left| \mathbf{E}_{b\mathbf{k}}\right| }=\frac{k_{b\bot }}{\sqrt{2}k_{b}}%
\left( \cos \alpha -i\sin \alpha \right) =\frac{k_{b\bot }}{\sqrt{2}k_{b}}%
e^{-i\alpha }  \notag \\
e_{b\mathbf{k},\Vert } &\equiv &\frac{E_{b\mathbf{k},z}}{\left| \mathbf{E}_{b%
\mathbf{k}}\right| }=\frac{k_{b\parallel }}{k_{b}}  \notag
\end{eqnarray}

\subsection{Determination of $\Theta _{\mathbf{k}}^{b,\text{LH}}$}

Lower Hybrid Current Drive (LHCD) results from momentum exchange from the LH
wave to the plasma through Landau damping (harmonic $n=0$). In this case,
the coefficient (\ref{thetak}) becomes 
\begin{eqnarray}
\Theta _{\mathbf{k}}^{b,0} &=&\frac{1}{\sqrt{2}}e_{b\mathbf{k},+}e^{-i\alpha
}J_{-1}\left( \frac{k_{\bot }v_{\bot }}{\Omega }\right) +\frac{1}{\sqrt{2}}%
e_{b\mathbf{k},-}e^{+i\alpha }J_{1}\left( \frac{k_{\bot }v_{\bot }}{\Omega }%
\right) +\frac{p_{\Vert }}{p_{\bot }}e_{b\mathbf{k},\Vert }J_{0}\left( \frac{%
k_{\bot }v_{\bot }}{\Omega }\right)  \notag \\
&=&\frac{1}{\sqrt{2}}\left( e_{b\mathbf{k},-}e^{+i\alpha }-e_{b\mathbf{k}%
,+}e^{-i\alpha }\right) J_{1}\left( \frac{k_{b\bot }v_{\bot }}{\Omega }%
\right) +\frac{p_{\Vert }}{p_{\bot }}e_{b\mathbf{k},\Vert }J_{0}\left( \frac{%
k_{b\bot }v_{\bot }}{\Omega }\right)
\end{eqnarray}

Using (\ref{pollh}), we see that the perpendicular components cancel and we
are left with 
\begin{equation}
\Theta _{\mathbf{k}}^{b,\text{LH}}=\frac{p_{\Vert }}{p_{\bot }}\frac{%
k_{b\parallel }}{k_{b}}J_{0}\left( \frac{k_{b\bot }v_{\bot }}{\Omega }\right)
\label{thetaklh}
\end{equation}

The argument of the Bessel function in (\ref{thetaklh}) is 
\begin{equation}
\frac{k_{b\bot }v_{\bot }}{\Omega }=N_{b\bot }\frac{v_{\bot }}{c}\frac{%
\omega _{b}}{\Omega }
\end{equation}
where an expression for $N_{\bot }$ is given by (\ref{nperplh1}). The
argument of the Bessel function becomes 
\begin{equation}
\left| \frac{k_{b\bot }v_{\bot }}{\Omega }\right| =\frac{1}{\sqrt{\omega
_{b}^{2}/\omega _{\text{LH}}^{2}-1}}\frac{\omega _{pe}}{\omega _{pi}}\frac{%
\omega _{b}}{\omega _{ce}}N_{b\parallel }\frac{p_{\bot }}{p_{Te}}\beta _{Te}
\end{equation}
and we see that for $\omega _{pe}\sim \omega _{ce}$ and $\omega _{b}\sim
2\omega _{\text{LH}}$, we get 
\begin{equation}
\left| \frac{k_{b\bot }v_{\bot }}{\Omega }\right| \sim \beta
_{Te}N_{b\parallel }\frac{p_{\bot }}{p_{Te}}
\end{equation}

Given that most electrons concerned with LHCD have $p_{\bot }\leq p_{Te}$,
and that for LHCD we have typically $N_{b\parallel }\sim 2$, we see that it
is reasonable to take the limit 
\begin{equation}
\left| \frac{k_{b\bot }v_{\bot }}{\Omega }\right| \ll 1
\end{equation}
as long as the plasma is not too relativistic ($\beta _{Te}\ll 1$). This
limit is consistent with the validity condition of the cold plasma
description. In this limit, valid for most LHCD scenarios, we have 
\begin{equation}
\Theta _{\mathbf{k}}^{b,\text{LH}}\simeq \frac{p_{\Vert }}{p_{\bot }}\frac{%
k_{b\parallel }}{k_{b}}  \label{thetaklh2}
\end{equation}

\subsection{Determination of $\mathbf{\Phi }_{b}^{\text{LH}}$}

The vector $\mathbf{\Phi }_{b}^{\text{LH}}$ describes the relation between
the energy flux and the electric field. Although LHW are almost
electrostatic, the dominant contribution to the power flux is the wave
Poynting flux, assuming that the wave frequency remains sensibly higher than
the LH frequency. The contribution of the Kinetic power flux is calculated
in Appendix \ref{kinetic}. In most LHCD scenarios, this contribution is not
more than a few percents of the total flux and can therefore be neglected.

In the cold plasma description, it is given by (\ref{phibres}). 
\begin{equation}
\mathbf{\Phi }_{b}=\left| e_{bT}\right| ^{2}\mathbf{N}_{b}-N_{b}\text{Re}%
\left( e_{bL}\mathbf{e}_{bT}^{\ast }\right)  \label{phiblh}
\end{equation}

LH waves are quasi-electrostatic, so that the longitudinal and transverse
components are given by (\ref{et}) 
\begin{eqnarray}
\left| e_{bL}\right| &\simeq &1 \\
\mathbf{e}_{bT} &=&\frac{1}{N_{b}^{2}}\mathbb{K}\cdot \mathbf{e}_{bL}
\end{eqnarray}
so that the first term in (\ref{phiblh}) can be neglected, and 
\begin{eqnarray}
\mathbf{\Phi }_{b}^{\text{LH}} &\simeq &-N_{b}\text{Re}\left( \varepsilon
_{bL}\frac{1}{N_{b}^{2}}\mathbb{K}^{\ast }\cdot \mathbf{e}_{bL}^{\ast
}\right)  \notag \\
&=&-\frac{\left| e_{bL}\right| ^{2}}{N_{b}}\text{Re}\left( \mathbb{K}^{\ast
}\cdot \widehat{N}\right)  \label{poyntinges}
\end{eqnarray}

Note from (\ref{poyntinges}) and (\ref{dispreles}) that the Poynting flux is
in the direction perpendicular to the wave vector: 
\begin{equation}
\mathbf{\Phi }_{b}^{\text{LH}}\cdot \widehat{N}=0
\end{equation}

We finally find, using (\ref{ktensor}), 
\begin{equation}
\mathbf{\Phi }_{b}^{\text{LH}}=-\frac{1}{N_{b}^{2}}\left( K_{\perp }\mathbf{N%
}_{b\perp }+K_{\Vert }N_{b\Vert }\widehat{z}\right)  \label{phiblh2}
\end{equation}

\subsection{LH Diffusion Coefficient}

The normalized bounce-averaged diffusion coefficient is given by (\ref
{diffcoefnorm3}) 
\begin{equation}
\overline{D}_{b,n}^{\text{RF}(0)}(p,\xi _{0})=\frac{\gamma p_{Te}}{p\left|
\xi _{0}\right| }\left| \Theta _{\mathbf{k},b}^{b,n}\left( \mathbf{k}%
_{b}\right) \right| ^{2}\overline{D}_{b,n,0}^{\text{RF},b}\Psi _{b}\frac{\xi
_{0}^{2}}{\xi _{b}^{2}}\frac{H_{\text{T}}\left( \xi _{0}\right) }{\lambda
\left( \xi _{0}\right) }\left[ \frac{1}{2}\sum_{\sigma =\pm 1}\right]
_{T}\delta \left( N_{\Vert b}-N_{\Vert \text{res},b}\right)  \label{dlh2}
\end{equation}
with 
\begin{equation}
\overline{D}_{b,n,0}^{\text{RF},b}=\frac{4\pi }{m_{e}\omega _{pe}^{2}\ln
\Lambda }\frac{\pi }{\omega _{b}}\frac{\Psi _{b}}{\Psi _{c}}\frac{1}{\left| 
\mathbf{\Phi }_{b}\cdot \widehat{n}_{b}\right| }\frac{P_{b,\text{inc}}}{4\pi
^{2}rR_{c}}
\end{equation}

Using (\ref{thetaklh2}), (\ref{zb}) gives 
\begin{equation}
\Theta _{\mathbf{k},b}^{b,\text{LH}}\simeq \frac{\xi _{b}}{\sqrt{\Psi
_{b}\left( 1-\xi _{0}^{2}\right) }}\frac{N_{b\parallel }}{N_{b}}
\end{equation}
so that (\ref{dlh2}) for LHCD\ can be written as 
\begin{equation}
\overline{D}_{b}^{\text{LH}(0)}(p,\xi _{0})=\frac{\gamma p_{Te}}{p\left| \xi
_{0}\right| }\frac{\xi _{0}^{2}}{\left( 1-\xi _{0}^{2}\right) }\frac{%
N_{b\parallel }^{2}}{N_{b}^{2}}\overline{D}_{b,0}^{\text{LH},b}\frac{H_{%
\text{T}}\left( \xi _{0}\right) }{\lambda \left( \xi _{0}\right) }\left[ 
\frac{1}{2}\sum_{\sigma =\pm 1}\right] _{T}\delta \left( N_{\Vert
b}-N_{\Vert \text{res},b}\right)
\end{equation}
where we have, using (\ref{phiblh2}) 
\begin{equation}
\overline{D}_{b,0}^{\text{LH},b}=\frac{4\pi }{m_{e}\omega _{pe}^{2}\ln
\Lambda }\frac{\pi }{\omega _{b}}\frac{N_{b}^{2}}{K_{\perp }N_{b\perp }}%
\frac{\Psi _{b}}{\Psi _{c}}\frac{1}{\left| \widehat{N}\cdot \widehat{n}%
_{b}\right| }\frac{P_{b,\text{inc}}}{4\pi ^{2}rR_{c}}
\end{equation}
and 
\begin{equation}
N_{\Vert \text{res},b}=\frac{1}{\beta _{Te}}\frac{\gamma p_{Te}}{p\xi _{b}}
\end{equation}

Using (\ref{nperplh1}), this becomes 
\begin{equation}
\overline{D}_{b}^{\text{LH}(0)}(p,\xi _{0})=\frac{\gamma p_{Te}}{p\left| \xi
_{0}\right| }\frac{\xi _{0}^{2}}{\left( 1-\xi _{0}^{2}\right) }\frac{%
K_{\perp }}{\left( -K_{\parallel }\right) }\overline{D}_{b,0}^{\text{LH},b}%
\frac{H_{\text{T}}\left( \xi _{0}\right) }{\lambda \left( \xi _{0}\right) }%
\left[ \frac{1}{2}\sum_{\sigma =\pm 1}\right] _{T}\delta \left( N_{\Vert
b}-N_{\Vert \text{res},b}\right)
\end{equation}
with 
\begin{equation}
\overline{D}_{b,0}^{\text{LH},b}=\frac{4\pi }{m_{e}\omega _{pe}^{2}\ln
\Lambda }\frac{\pi }{\omega _{b}}\frac{N_{b\parallel }}{K_{\perp }}\sqrt{%
\frac{-K_{\parallel }}{K_{\perp }}}\frac{\Psi _{b}}{\Psi _{c}}\frac{1}{%
\left| \widehat{N}\cdot \widehat{n}_{b}\right| }\frac{P_{b,\text{inc}}}{4\pi
^{2}rR_{c}}
\end{equation}
or, expliciting the dielectric tensor elements, 
\begin{equation}
\overline{D}_{b}^{\text{LH}(0)}(p,\xi _{0})=\frac{\gamma p_{Te}}{p\left| \xi
_{0}\right| }\frac{\xi _{0}^{2}}{\left( 1-\xi _{0}^{2}\right) }\frac{\omega
_{pi}^{2}}{\omega _{pe}^{2}}\left( \omega ^{2}/\omega _{\text{LH}%
}^{2}-1\right) \overline{D}_{b,0}^{\text{LH},b}\frac{H_{\text{T}}\left( \xi
_{0}\right) }{\lambda \left( \xi _{0}\right) }\left[ \frac{1}{2}\sum_{\sigma
=\pm 1}\right] _{T}\delta \left( N_{\Vert b}-N_{\Vert \text{res},b}\right)
\end{equation}
with 
\begin{equation}
\overline{D}_{b,0}^{\text{LH},b}=\frac{4\pi ^{2}}{m_{e}\omega _{pe}\omega
^{2}\ln \Lambda }\frac{N_{b\parallel }}{\left[ \left( 1+\omega
_{pe}^{2}/\omega _{ce}^{2}\right) \left( 1-\omega _{\text{LH}}^{2}/\omega
^{2}\right) \right] ^{3/2}}\frac{\Psi _{b}}{\Psi _{c}}\frac{1}{\left| 
\widehat{N}\cdot \widehat{n}_{b}\right| }\frac{P_{b,\text{inc}}}{4\pi
^{2}rR_{c}}
\end{equation}

At this stage, the only wave characteristics needed to calculate the
diffusion coefficient are $\theta _{b}$, $N_{\Vert b}$ and $P_{b,\text{inc}}$
for each ray.

\subsubsection{Square power spectrum in $k_{\Vert }$}

\bigskip The power spectrum for LH waves is commonly taken as being constant
in $N_{\Vert }$ between two limits $N_{\Vert \min }$ and $N_{\Vert \max }$\
that are respectively given by the accessibility condition and the condition
for strong linear Laudau damping. Strictly speaking, we should decompose
this spectrum in a sum of rays with narrow $N_{\Vert }$ extension. If we
make the approximation that the characteristics of each ray are quite
similar, which is a good approximation only when 
\begin{equation}
\Delta N_{\Vert }=N_{\Vert \max }-N_{\Vert \min }\ll N_{\Vert }
\end{equation}
then it is possible to consider only one ray, by transforming 
\begin{equation}
\delta \left( N_{\Vert b}-N_{\Vert \text{res},b}\right) \rightarrow \frac{1}{%
\Delta N_{\Vert }}H\left( N_{\Vert \text{res},b}-N_{\Vert \min }\right)
H\left( N_{\Vert \max }-N_{\Vert \text{res},b}\right)
\end{equation}

The LH diffusion coefficient then becomes 
\begin{equation}
\overline{D}_{b}^{\text{LH}(0)}(p,\xi _{0})=\frac{\gamma p_{Te}}{p\left| \xi
_{0}\right| }\frac{\xi _{0}^{2}}{\left( 1-\xi _{0}^{2}\right) }\overline{D}%
_{b,0}^{\text{LH},b}\frac{N_{\parallel ,0}^{2}}{N_{0}^{2}}\frac{H_{\text{T}%
}\left( \xi _{0}\right) }{\lambda \left( \xi _{0}\right) }\frac{1}{\Delta
N_{\Vert }}\left[ \frac{1}{2}\sum_{\sigma =\pm 1}\right] _{T}H\left(
N_{\Vert \text{res},b}-N_{\Vert \min }\right) H\left( N_{\Vert \max
}-N_{\Vert \text{res},b}\right)  \label{dlh}
\end{equation}
with 
\begin{equation}
\overline{D}_{b,0}^{\text{LH},b}=\frac{4\pi ^{2}}{m_{e}\omega
_{pe}^{2}\omega \ln \Lambda }\frac{N_{0}^{2}}{K_{\perp }N_{\perp ,0}}\frac{%
\Psi _{b}}{\Psi _{c}}\frac{1}{\left| \widehat{N}\cdot \widehat{n}_{b}\right| 
}\frac{P_{\text{inc}}}{4\pi ^{2}rR_{c}}  \label{d0lh}
\end{equation}

\subsubsection{Ignoring trapped particle effects.}

Trapped particle effects can be readily ignored when $\epsilon \rightarrow 0$
(central CD) where the trapped fraction vanishes. In this case, 
\begin{eqnarray}
\Psi _{b},\Psi _{c} &\rightarrow &1 \\
\lambda \left( \xi _{0}\right) &\rightarrow &1 \\
H_{\text{T}}\left( \xi _{0}\right) &\rightarrow &1 \\
\xi _{b} &\rightarrow &\xi _{0}
\end{eqnarray}
and we find 
\begin{equation}
\overline{D}_{b}^{\text{LH}}(p,\xi _{0})=\frac{\gamma p_{Te}}{p\left| \xi
_{0}\right| }\frac{\xi _{0}^{2}}{\left( 1-\xi _{0}^{2}\right) }\overline{D}%
_{b,0}^{\text{LH},b}\frac{N_{\parallel ,0}^{2}}{N_{0}^{2}}\frac{1}{\Delta
N_{\Vert }}H\left( N_{\Vert \text{res}}-N_{\Vert \min }\right) H\left(
N_{\Vert \max }-N_{\Vert \text{res}}\right)  \label{dlh_notrap}
\end{equation}
where 
\begin{equation}
N_{\Vert \text{res}}=\frac{1}{\beta _{Te}}\frac{\gamma p_{Te}}{p\xi _{0}}
\end{equation}
and 
\begin{equation}
\overline{D}_{b,0}^{\text{LH},b}=\frac{4\pi ^{2}}{m_{e}\omega
_{pe}^{2}\omega \ln \Lambda }\frac{N_{0}^{2}}{K_{\perp }N_{\perp ,0}}\frac{1%
}{\left| \widehat{N}\cdot \widehat{n}_{b}\right| }\frac{P_{\text{inc}}}{4\pi
^{2}rR_{c}}
\end{equation}

\subsubsection{Common form of the LH coefficient in the Litterature}

In order to retrieve a more common form for the LH QL Operator, we can
rewritte (\ref{dlh_notrap}) as 
\begin{equation}
\overline{D}_{b}^{\text{LH}}(p,\xi _{0})=\frac{\gamma p_{Te}}{p\left| \xi
_{0}\right| }\frac{\xi _{0}^{2}}{\left( 1-\xi _{0}^{2}\right) }\overline{D}%
_{0}^{\text{LH}}H\left( N_{\Vert \text{res}}-N_{\Vert \min }\right) H\left(
N_{\Vert \max }-N_{\Vert \text{res}}\right)
\end{equation}
where we define 
\begin{eqnarray}
\overline{D}_{0}^{\text{LH}} &=&\overline{D}_{b,0}^{\text{LH},b}\frac{%
N_{\parallel ,0}^{2}}{N_{0}^{2}}\frac{1}{\Delta N_{\Vert }} \\
&=&\frac{4\pi ^{2}}{m_{e}\omega _{pe}^{2}\omega \ln \Lambda }\frac{1}{\sqrt{%
-K_{\parallel }K_{\perp }}}\frac{N_{\parallel ,0}}{\Delta N_{\Vert }}\frac{1%
}{\left| \widehat{N}\cdot \widehat{n}_{b}\right| }\frac{P_{\text{inc}}}{4\pi
^{2}rR_{c}}  \notag \\
&=&\frac{4\pi ^{2}}{m_{e}\omega _{pe}^{3}\ln \Lambda }\frac{1}{\sqrt{%
1+\omega _{pe}^{2}/\omega _{ce}^{2}}\sqrt{1-\omega _{\text{LH}}^{2}/\omega
^{2}}}\frac{N_{\parallel ,0}}{\Delta N_{\Vert }}\frac{1}{\left| \widehat{N}%
\cdot \widehat{n}_{b}\right| }\frac{P_{\text{inc}}}{4\pi ^{2}rR_{c}}
\label{dlh1}
\end{eqnarray}

A familiar expression for the LH QL operator is obtained in cylindrical
geometry. From (\ref{dcyl}) and (\ref{divcyl}) we see that for LHCD, the QL\
operator can be rewritten as 
\begin{eqnarray}
Q^{\text{LH}}(f) &=&\frac{\partial }{\partial p_{\Vert }}D_{\Vert \Vert }^{%
\text{RF}}\frac{\partial f}{\partial p_{\Vert }}  \notag \\
&=&\frac{\partial }{\partial p_{\Vert }}D_{0}^{\text{LH}}\frac{v_{Te}}{%
\left| v_{\Vert }\right| }H\left( \frac{c}{v_{\Vert }}-N_{\Vert \min
}\right) H\left( N_{\Vert \max }-\frac{c}{v_{\Vert }}\right) \frac{\partial f%
}{\partial p_{\Vert }}
\end{eqnarray}

\subsubsection{Simple Approximate Form of the LH QL Operator.}

The factor $v_{Te}/\left| v_{\Vert }\right| $ has often been neglected in
the litterature, which is an acceptable approximation when $D_{0}^{\text{LH}%
}\gg \nu _{e}p_{Te}^{2}$ and $\Delta N_{\Vert }\ll N_{\Vert \min },N_{\Vert
\max }$. In this case, the LH Operator reduces to 
\begin{equation}
Q^{\text{LH}}(f)\simeq \left\{ 
\begin{array}{cc}
\frac{\partial }{\partial p_{\Vert }}D_{0}^{\text{LH}}\frac{\partial f}{%
\partial p_{\Vert }} & \text{for }v_{1}<v_{\Vert }/v_{Te}<v_{2} \\ 
0 & \text{otherwise}
\end{array}
\right.
\end{equation}
with 
\begin{eqnarray}
v_{1} &=&\frac{1}{\beta _{Te}N_{\Vert \max }} \\
v_{2} &=&\frac{1}{\beta _{Te}N_{\Vert \min }}
\end{eqnarray}

However, in a Fokker-Planck code, it is better to keep the exact formula (%
\ref{dlh}). Note that $N_{\parallel ,0}$ is not really defined in (\ref{dlh1}%
). We can take 
\begin{equation}
N_{\parallel ,0}=\frac{N_{\Vert \min }+N_{\Vert \max }}{2}
\end{equation}
which leads approximately to 
\begin{equation}
\frac{N_{\parallel }}{\Delta N_{\parallel }}\simeq \frac{2v_{1}v_{2}}{%
v_{2}^{2}-v_{1}^{2}}
\end{equation}

\section{Diffusion coefficient for Electron Cyclotron Current Drive}

The Cold Plasma description is usually a good approximation to determine the
EC Wave spectrum properties and polarization, as long as we stay away from
the upper-hybrid resonance, where mode-conversion to EBW occurs. However,
even in the cold plasma model, the polarizations are usually mixed for
oblique propagation, and no simple analytical formulation is possible.

The prescription is to decompose the $N_{\parallel }$ spectrum into many
rays $b$, each of which has a narrow spectrum around $N_{\parallel b}$ and
therefore can be associated to a plane wave and treated as described in
Section \ref{beam}. Then, one needs to solve the dispersion relation to find 
$N_{\bot b}$ as a function of $N_{\parallel b}$, and solve the wave equation
to find the polarization components, as described in Appendix \ref{cold}.
The quantities $\Theta _{\mathbf{k},b}^{b,n}\left( \mathbf{k}_{b}\right) $\
and\ $\mathbf{\Phi }_{b}$ can then be evaluated to determine the diffusion
coefficient (\ref{diffcoefnorm3}).

Still, it is possible to find limit cases (small FLR effects, small $%
N_{\Vert }$) for which a simple analytical derivation is possible. After
simplifying the problem by introducing the finite size beam, some limit
cases are investigated for the X-2 and the O-1 modes.

\subsection{Finite size beam}

When the total $k_{\parallel }$ spectrum is rather narrowly centered around
some value $k_{\parallel ,0}$, that is 
\begin{equation}
\frac{\Delta k_{\parallel }}{k_{\parallel ,0}}\ll 1  \label{narrow}
\end{equation}
then the characteristics of the rays (spectrum and polarization) can be very
similar, and decomposing may become unnecessary. This can be the case
particularly for EC beams.

In this case, we can consider only one ray, treated as described in Section 
\ref{beam}, and replace the delta function by the actual spectrum, for
example, a Gaussian spectrum 
\begin{equation}
\delta \left( N_{\Vert b}-N_{\Vert \text{res}}\right) \rightarrow \frac{1}{%
\sqrt{\pi }\Delta N_{\Vert }}\exp \left[ \frac{-\left( N_{\Vert \text{res}%
}-N_{\Vert b}\right) ^{2}}{\Delta N_{\Vert }^{2}}\right]
\end{equation}

By the properties of Fourier transform, the width $\Delta N_{\Vert }$ is
related to the ray waist $d$, 
\begin{equation}
\Delta N_{\Vert }=\frac{N_{\bot b}}{N_{b}}\frac{c}{\omega _{b}d}
\end{equation}
and we can verify that for waist of a few centimeters, the condition (\ref
{narrow}) is satisfied. For such beams, using (\ref{diffcoefnorm3}), the EC
diffusion coefficient becomes 
\begin{equation}
\overline{D}_{b,n}^{\text{EC}(0)}(p,\xi _{0})=\frac{\gamma p_{Te}}{p\left|
\xi _{0}\right| }\left| \Theta _{\mathbf{k},b}^{b,n}\left( \mathbf{k}%
_{b}\right) \right| ^{2}\overline{D}_{b,n,0}^{\text{EC},b}\Psi _{b}\frac{\xi
_{0}^{2}}{\xi _{b}^{2}}\frac{H_{\text{T}}\left( \xi _{0}\right) }{\lambda
\left( \xi _{0}\right) }\frac{1}{\sqrt{\pi }\Delta N_{\Vert }}\left[ \frac{1%
}{2}\sum_{\sigma =\pm 1}\right] _{T}\exp \left[ \frac{-\left( N_{\Vert \text{%
res},b}-N_{\Vert b}\right) ^{2}}{\Delta N_{\Vert }^{2}}\right]  \label{dec}
\end{equation}
with 
\begin{equation}
\overline{D}_{b,n,0}^{\text{EC},b}=\frac{4\pi ^{2}}{m_{e}\omega
_{pe}^{2}\omega _{b}\ln \Lambda }\frac{\Psi _{b}}{\Psi _{c}}\frac{1}{\left| 
\mathbf{\Phi }_{b}\cdot \widehat{n}_{b}\right| }\frac{P_{b,\text{inc}}}{4\pi
^{2}rR_{c}}
\end{equation}

For realistic calculations, the polarization coefficient $\Theta _{\mathbf{k}%
,b}^{b,n}\left( \mathbf{k}_{b}\right) $ must be calculated exactly (See
Appendix \ref{cold}). However, to get some physical insight, and also for
the sake of comparison with other works, it can be useful to consider some
limit cases.

\subsection{Polarization: Limit of small FLR effects}

Electron Cyclotron Current Drive (LHCD) results from momentum exchange from
the LH wave to the plasma through Cyclotron damping at some harmonic $n$.
For a beam $b$, considering a given harmonic $n$, the coefficient (\ref
{thetak}) is 
\begin{equation}
\Theta _{\mathbf{k}}^{b,n}=\frac{1}{\sqrt{2}}e_{b\mathbf{k},+}e^{-i\alpha
}J_{n-1}\left( z\right) +\frac{1}{\sqrt{2}}e_{b\mathbf{k},-}e^{+i\alpha
}J_{n+1}\left( z\right) +\frac{p_{\Vert }}{p_{\bot }}e_{b\mathbf{k},\Vert
}J_{n}\left( z\right)
\end{equation}

The resonance condition 
\begin{equation}
\omega -k_{\Vert }v_{\Vert }-n\Omega
\end{equation}
implies that resonant particles are those who verify 
\begin{equation}
\frac{v_{\Vert }}{c}=\frac{1}{N_{\Vert }}\left( 1-\frac{n\Omega }{\omega }%
\right)
\end{equation}

Given that generally $N_{\Vert }\lesssim 1$, there are resonant electrons
(with $v_{\Vert }<c$) only if 
\begin{equation}
\frac{n\Omega }{\omega }\sim 1  \label{ysim}
\end{equation}

Since 
\begin{equation}
\Omega =\frac{q_{e}B}{\gamma m_{e}}=-\frac{\omega _{ce}}{\gamma }<0
\end{equation}
the only harmonics to be considered for electrons are for 
\begin{equation}
n\leq -1
\end{equation}

Applying the substitution 
\begin{equation}
n^{\prime }=-n\geq 1
\end{equation}
and renaming $n^{\prime }\rightarrow n$, we get 
\begin{eqnarray}
\Theta _{\mathbf{k}}^{b,-n} &=&\frac{1}{\sqrt{2}}e_{b\mathbf{k}%
,+}e^{-i\alpha }J_{-n-1}\left( z\right) +\frac{1}{\sqrt{2}}e_{b\mathbf{k}%
,-}e^{+i\alpha }J_{-n+1}\left( z\right) +\frac{p_{\Vert }}{p_{\bot }}e_{b%
\mathbf{k},\Vert }J_{-n}\left( z\right) \\
&=&\frac{1}{\sqrt{2}}e_{b\mathbf{k},+}e^{-i\alpha }\left( -1\right)
^{n+1}J_{n+1}\left( z\right) +\frac{1}{\sqrt{2}}e_{b\mathbf{k},-}e^{+i\alpha
}\left( -1\right) ^{n-1}J_{n-1}\left( z\right) +\frac{p_{\Vert }}{p_{\bot }}%
e_{b\mathbf{k},\Vert }\left( -1\right) ^{n}J_{n}\left( z\right)
\end{eqnarray}

We have 
\begin{equation}
z=\frac{k_{\bot }v_{\bot }}{\Omega }=-N_{\bot }\frac{p_{\bot }}{m_{e}c}\frac{%
\omega _{b}}{\omega _{ce}}
\end{equation}
Using (\ref{ysim}), we get an estimation for $\left| z\right| $%
\begin{equation}
\left| z\right| \sim nN_{\bot }\frac{p_{\bot }}{p_{Te}}\beta _{Te}
\end{equation}
Typically, for ECCD in tokamaks, we have 
\begin{equation}
N_{\bot }\lesssim 1,\;\beta _{Te}\lesssim 0.1
\end{equation}
and most electrons of concern have 
\begin{equation}
\frac{p_{\bot }}{p_{Te}}\leq 3
\end{equation}
so that for the first and second harmonics, the condition $z\ll 1$ is
satisfied. In other words, the Larmor radius remains small compared to the
perpendicular wavelength. This condition is consistent with the conditions
of the cold plasma description of the EC waves. In this case, for $n\geq 0$,
the following approximate expression can be used 
\begin{equation}
J_{n}\left( z\right) \simeq \frac{1}{n!}\left( \frac{z}{2}\right) ^{n}
\end{equation}
and we have 
\begin{equation}
\Theta _{\mathbf{k}}^{b,-n}\simeq \frac{e_{b\mathbf{k},+}}{\sqrt{2}}%
e^{-i\alpha }\left( -1\right) ^{n+1}\frac{1}{\left( n+1\right) !}\left( 
\frac{z}{2}\right) ^{n+1}+\frac{e_{b\mathbf{k},-}}{\sqrt{2}}e^{+i\alpha
}\left( -1\right) ^{n-1}\frac{1}{\left( n-1\right) !}\left( \frac{z}{2}%
\right) ^{n-1}+\frac{p_{\Vert }}{p_{\bot }}e_{b\mathbf{k},\Vert }\frac{1}{n!}%
\left( \frac{z}{2}\right) ^{n}  \label{thetaklowz}
\end{equation}

\subsection{X-Mode, $n=2$}

\subsubsection{Polarization}

In the X-Mode polarization, the harmonic $n=1$ is particular in the sense
that the polarization component $e_{b\mathbf{k},-}$ vanishes. In this case,
absorption is significantly reduced, because $e_{b\mathbf{k},-}$ is related
to the lowest power of $z$. Moreover, this harmonic can only be reached from
the high field side, because of the right-hand cut-off. We therefore
consider the second harmonic $n=2$. For this harmonic, $e_{b\mathbf{k},-}$
is usually the largest component of the polarization, and, since it is
associated with the lowest power of $z$ with $z\ll 1$, the corresponding
term is the dominant term in (\ref{thetaklowz}), which reduces to 
\begin{equation}
\Theta _{\mathbf{k}}^{b,\text{X2}}\simeq N_{\bot }\frac{p_{\bot }}{m_{e}c}%
\frac{\omega _{b}}{2\omega _{ce}}\frac{e_{b\mathbf{k},-}}{\sqrt{2}}%
e^{+i\alpha }
\end{equation}

Then, using (\ref{zb}-\ref{zb2}), we find 
\begin{equation}
\Theta _{\mathbf{k},b}^{b,\text{X2}}\simeq N_{\bot b}\frac{p}{m_{e}c}\frac{%
\sqrt{1-\xi _{0}^{2}}}{\sqrt{\Psi _{b}}}\frac{\omega _{b}}{2\omega _{ce,0}}%
\frac{e_{b\mathbf{k},-}}{\sqrt{2}}e^{+i\alpha }
\end{equation}

Then, the diffusion coefficient (\ref{dec}) becomes 
\begin{eqnarray}
\overline{D}_{b}^{\text{EC-X2}(0)}(p,\xi _{0}) &=&\frac{\gamma p_{Te}}{%
p\left| \xi _{0}\right| }\left( \frac{p}{p_{Te}}\right) ^{2}\left( 1-\xi
_{0}^{2}\right) \beta _{Te}^{2}N_{\bot b}^{2}\left( \frac{\omega _{b}}{%
2\omega _{ce,0}}\right) ^{2}\frac{\left| e_{b,-}\right| ^{2}}{2}\overline{D}%
_{0}^{\text{EC-X2}(0)}\frac{\xi _{0}^{2}}{\xi _{b}^{2}}\frac{H_{\text{T}%
}\left( \xi _{0}\right) }{\lambda \left( \xi _{0}\right) }  \label{decx2} \\
&&\frac{1}{\sqrt{\pi }\Delta N_{\Vert }}\left[ \frac{1}{2}\sum_{\sigma =\pm
1}\right] _{T}\exp \left[ \frac{-\left( N_{\Vert \text{res},b}-N_{\Vert
b}\right) ^{2}}{\Delta N_{\Vert }^{2}}\right]  \notag
\end{eqnarray}
with 
\begin{equation}
\overline{D}_{0}^{\text{EC-X2}(0)}=\frac{4\pi }{m_{e}\omega _{pe}^{2}\ln
\Lambda }\frac{\pi }{\omega _{b}}\frac{\Psi _{b}}{\Psi _{c}}\frac{1}{\left| 
\mathbf{\Phi }_{b}\cdot \widehat{n}_{b}\right| }\frac{P_{b,\text{inc}}}{4\pi
^{2}rR_{c}}
\end{equation}

\subsubsection{Determination of $\mathbf{\Phi }_{b}^{\text{EC-X2}(0)}$ in
the low density limit.}

The vector $\mathbf{\Phi }_{b}$ describes the relation between the energy
flux and the electric field. In the cold plasma description, it is given by (%
\ref{phibres}). 
\begin{equation}
\mathbf{\Phi }_{b}=N_{b}\left[ \left| e_{bT}\right| ^{2}\widehat{N}_{b}-%
\text{Re}\left( e_{bL}\mathbf{e}_{bT}^{\ast }\right) \right]
\end{equation}

In general, both terms must be kept. In the limit of low density 
\begin{equation}
\omega _{pe}\ll \omega _{ce}
\end{equation}
the longitudinal component of the polarization vanishes and 
\begin{eqnarray}
\left| e_{bL}\right| &\ll &\left| e_{bT}\right| \\
\left| e_{bT}\right| &\simeq &1 \\
N_{b} &\simeq &1
\end{eqnarray}
so that 
\begin{equation}
\mathbf{\Phi }_{b}^{\omega _{pe}\ll \omega _{ce}}\simeq \widehat{N}_{b}
\end{equation}

We then find 
\begin{equation}
\overline{D}_{0\text{,}\omega _{pe}\ll \omega _{ce}\text{ }}^{\text{EC-X2}%
(0)}=\frac{4\pi }{m_{e}\omega _{pe}^{2}\ln \Lambda }\frac{\pi }{\omega _{b}}%
\frac{\Psi _{b}}{\Psi _{c}}\frac{1}{\left| \widehat{N}_{b}\cdot \widehat{n}%
_{b}\right| }\frac{P_{b,\text{inc}}}{4\pi ^{2}rR_{c}}
\end{equation}

\subsubsection{Common form of the EC-X2 coefficient in the Litterature}

In order to retrieve a more common form for the LH EC-X2 Operator, we can
rewritte (\ref{decx2}) as 
\begin{equation}
\overline{D}_{b}^{\text{EC-X2}(0)}(p,\xi _{0})=\frac{\gamma p_{Te}}{p\left|
\xi _{0}\right| }\left( \frac{p}{p_{Te}}\right) ^{2}\left( 1-\xi
_{0}^{2}\right) \overline{D}_{0,\text{redef}}^{\text{EC-X2}(0)}\frac{\xi
_{0}^{2}}{\xi _{b}^{2}}\frac{H_{\text{T}}\left( \xi _{0}\right) }{\lambda
\left( \xi _{0}\right) }\left[ \frac{1}{2}\sum_{\sigma =\pm 1}\right]
_{T}\exp \left[ \frac{-\left( N_{\Vert \text{res},b}-N_{\Vert b}\right) ^{2}%
}{\Delta N_{\Vert }^{2}}\right]  \notag
\end{equation}
with 
\begin{eqnarray}
\overline{D}_{0,\text{redef}}^{\text{EC-X2}(0)} &=&\frac{1}{\sqrt{\pi }%
\Delta N_{\Vert }}\beta _{Te}^{2}N_{\bot b}^{2}\left( \frac{\omega _{b}}{%
2\omega _{ce,0}}\right) ^{2}\frac{\left| e_{b,-}\right| ^{2}}{2}\overline{D}%
_{0}^{\text{EC-X2}(0)} \\
&=&\frac{4\pi }{m_{e}\omega _{pe}^{2}\ln \Lambda }\frac{\sqrt{\pi }}{\omega
_{b}\Delta N_{\Vert }}\beta _{Te}^{2}N_{\bot b}^{2}\left( \frac{\omega _{b}}{%
2\omega _{ce,0}}\right) ^{2}\frac{\Psi _{b}}{\Psi _{c}}\frac{\left|
e_{b,-}\right| ^{2}}{2\left| \mathbf{\Phi }_{b}\cdot \widehat{n}_{b}\right| }%
\frac{P_{b,\text{inc}}}{4\pi ^{2}rR_{c}}
\end{eqnarray}

\subsection{O-Mode, $n=1,N_{\Vert }\ll 1$}

\subsubsection{Polarization}

In the O-Mode polarization, the harmonic $n=1$ usually gives the largest
absorption, and is accessible from the LFS, provided that $\omega
_{pe}<\omega _{ce}$. For this harmonic, (\ref{thetaklowz}) becomes 
\begin{equation}
\Theta _{\mathbf{k}}^{b,\text{O1}}\simeq \frac{e_{b\mathbf{k},+}}{\sqrt{2}}%
e^{-i\alpha }\frac{1}{2}\left( \frac{z}{2}\right) ^{2}+\frac{e_{b\mathbf{k}%
,-}}{\sqrt{2}}e^{+i\alpha }+\frac{p_{\Vert }}{p_{\bot }}e_{b\mathbf{k},\Vert
}\frac{z}{2}
\end{equation}

Here the hierarchy between polarization components is more complicated. For
small toroidal angle ($N_{\Vert }\ll 1$), the parallel component $e_{b%
\mathbf{k},\Vert }$ dominates largely and, if the temperature is not too
low, we can have 
\begin{equation}
\left| \frac{p_{\Vert }}{p_{\bot }}e_{b\mathbf{k},\Vert }\frac{z}{2}\right|
\gg \frac{\left| e_{b\mathbf{k},-}\right| }{\sqrt{2}},\;\frac{\left| e_{b%
\mathbf{k},+}\right| }{\sqrt{2}}\frac{1}{2}\left( \frac{z}{2}\right) ^{2}
\end{equation}

In this case, we can take 
\begin{equation}
\Theta _{\mathbf{k}}^{b,\text{O1,}N_{\Vert }\ll 1}\simeq e_{b\mathbf{k}%
,\Vert }N_{\bot }\frac{p_{\Vert }}{m_{e}c}\frac{\omega _{b}}{2\omega _{ce}}
\end{equation}
so that, using (\ref{zb}-\ref{zb2}), we find 
\begin{equation}
\Theta _{\mathbf{k},b}^{b,\text{O1,}N_{\Vert }\ll 1}\simeq N_{\bot b}\frac{p%
}{m_{e}c}\frac{\xi _{b}}{\Psi _{b}}\frac{\omega _{b}}{2\omega _{ce,0}}e_{b%
\mathbf{k},\Vert }
\end{equation}

Then, the diffusion coefficient (\ref{dec}) can be expressed as 
\begin{eqnarray}
\overline{D}_{b,N_{\Vert }\ll 1}^{\text{EC-O1}(0)}(p,\xi _{0}) &=&\frac{%
\gamma p_{Te}}{p\left| \xi _{0}\right| }\left( \frac{p}{p_{Te}}\right)
^{2}N_{\bot b}^{2}\beta _{Te}^{2}\left( \frac{\omega _{b}}{2\omega _{ce,0}}%
\right) ^{2}\left| e_{b\mathbf{k},\Vert }\right| ^{2}\overline{D}%
_{0,N_{\Vert }\ll 1}^{\text{EC-O1}(0)}\frac{\xi _{0}^{2}}{\Psi _{b}}\frac{H_{%
\text{T}}\left( \xi _{0}\right) }{\lambda \left( \xi _{0}\right) }
\label{deco1} \\
&&\frac{1}{\sqrt{\pi }\Delta N_{\Vert }}\left[ \frac{1}{2}\sum_{\sigma =\pm
1}\right] _{T}\exp \left[ \frac{-\left( N_{\Vert \text{res},b}-N_{\Vert
b}\right) ^{2}}{\Delta N_{\Vert }^{2}}\right]  \notag
\end{eqnarray}
with 
\begin{equation}
\overline{D}_{0,N_{\Vert }\ll 1}^{\text{EC-O1}(0)}=\frac{4\pi }{m_{e}\omega
_{pe}^{2}\ln \Lambda }\frac{\pi }{\omega _{b}}\frac{\Psi _{b}}{\Psi _{c}}%
\frac{1}{\left| \mathbf{\Phi }_{b}\cdot \widehat{n}_{b}\right| }\frac{P_{b,%
\text{inc}}}{4\pi ^{2}rR_{c}}
\end{equation}

\subsubsection{Determination of $\mathbf{\Phi }_{b}^{\text{EC-O}(0)}$}

The vector $\mathbf{\Phi }_{b}$ describes the relation between the energy
flux and the electric field. In the cold plasma description, it is given by (%
\ref{phibres}). 
\begin{equation}
\mathbf{\Phi }_{b}=N_{b}\left[ \left| e_{bT}\right| ^{2}\widehat{N}_{b}-%
\text{Re}\left( e_{bL}\mathbf{e}_{bT}^{\ast }\right) \right]
\end{equation}

For small toroidal angle ($N_{\Vert }\ll 1$), we have 
\begin{equation}
\mathbf{e}_{b}\simeq e_{b\mathbf{k},\Vert }\widehat{z}
\end{equation}
and, since $N_{\Vert }\ll 1$, the longitudinal component of $\varepsilon _{b%
\mathbf{k},\Vert }$ is very small, so that 
\begin{equation}
\mathbf{\Phi }_{b}^{\text{EC-O}(0)}\simeq \mathbf{N}_{b}
\end{equation}

We then find 
\begin{equation}
\overline{D}_{0,N_{\Vert }\ll 1}^{\text{EC-O1}(0)}=\frac{4\pi }{m_{e}\omega
_{pe}^{2}\ln \Lambda }\frac{\pi }{\omega _{b}}\frac{1}{N_{b}}\frac{1}{\left| 
\widehat{N}_{b}\cdot \widehat{n}_{b}\right| }\frac{P_{b,\text{inc}}}{4\pi
^{2}rR_{c}}
\end{equation}

\subsubsection{Common form of the EC-O1 coefficient in the Litterature}

In order to retrieve a more common form for the LH EC-O1 Operator, we can
rewritte (\ref{deco1}) as 
\begin{equation}
\overline{D}_{b,N_{\Vert }\ll 1}^{\text{EC-O1}(0)}(p,\xi _{0})=\frac{\gamma
p_{Te}}{p\left| \xi _{0}\right| }\left( \frac{p}{p_{Te}}\right) ^{2}%
\overline{D}_{0,N_{\Vert }\ll 1,\text{redef}}^{\text{EC-O1}(0)}\frac{\xi
_{0}^{2}}{\Psi _{b}}\frac{H_{\text{T}}\left( \xi _{0}\right) }{\lambda
\left( \xi _{0}\right) }\left[ \frac{1}{2}\sum_{\sigma =\pm 1}\right]
_{T}\exp \left[ \frac{-\left( N_{\Vert \text{res},b}-N_{\Vert b}\right) ^{2}%
}{\Delta N_{\Vert }^{2}}\right]  \notag
\end{equation}
with 
\begin{eqnarray}
\overline{D}_{0,N_{\Vert }\ll 1,\text{redef}}^{\text{EC-O1}(0)} &=&\frac{1}{%
\sqrt{\pi }\Delta N_{\Vert }}N_{\bot b}^{2}\beta _{Te}^{2}\left( \frac{%
\omega _{b}}{2\omega _{ce,0}}\right) ^{2}\left| e_{b\mathbf{k},\Vert
}\right| ^{2}\overline{D}_{0,N_{\Vert }\ll 1}^{\text{EC-O1}(0)} \\
&=&\frac{4\pi }{m_{e}\omega _{pe}^{2}\ln \Lambda }\frac{\sqrt{\pi }}{\omega
_{b}\Delta N_{\Vert }}\frac{N_{\bot b}^{2}}{N_{b}}\beta _{Te}^{2}\left( 
\frac{\omega _{b}}{2\omega _{ce,0}}\right) ^{2}\frac{\Psi _{b}}{\Psi _{c}}%
\frac{1}{\left| \widehat{N}_{b}\cdot \widehat{n}_{b}\right| }\frac{P_{b,%
\text{inc}}}{4\pi ^{2}rR_{c}}
\end{eqnarray}

\appendix

\section{Transformation of the QL Operator\label{fluxexp}}

The transformation from cylindrical to spherical momentum derivatives is
given by 
\begin{eqnarray}
\frac{\partial }{\partial p_{\Vert }} &=&\frac{\partial p}{\partial p_{\Vert
}}\frac{\partial }{\partial p}+\frac{\partial \xi }{\partial p_{\Vert }}%
\frac{\partial }{\partial \xi }  \label{dppar} \\
\frac{\partial }{\partial p_{\bot }} &=&\frac{\partial p}{\partial p_{\bot }}%
\frac{\partial }{\partial p}+\frac{\partial \xi }{\partial p_{\bot }}\frac{%
\partial }{\partial \xi }  \label{dpperp}
\end{eqnarray}

We have 
\begin{eqnarray}
p &=&\sqrt{p_{\Vert }^{2}+p_{\bot }^{2}} \\
\xi &=&\frac{p_{\Vert }}{\sqrt{p_{\Vert }^{2}+p_{\bot }^{2}}}
\end{eqnarray}
so that 
\begin{eqnarray}
\frac{\partial p}{\partial p_{\Vert }} &=&\xi \\
\frac{\partial p}{\partial p_{\bot }} &=&\sqrt{1-\xi ^{2}} \\
\frac{\partial \xi }{\partial p_{\Vert }} &=&\frac{1-\xi ^{2}}{p} \\
\frac{\partial \xi }{\partial p_{\bot }} &=&-\frac{\xi \sqrt{1-\xi ^{2}}}{p}
\end{eqnarray}

We rewritte (\ref{ppar}-\ref{pperp}) in cylindrical coordinates using (\ref
{dppar}-\ref{dpperp}) 
\begin{eqnarray}
P_{\Vert } &=&\xi \frac{\partial }{\partial p}+\frac{1}{p}\left( 1-\xi ^{2}-%
\frac{n\Omega }{\omega }\right) \frac{\partial }{\partial \xi }
\label{ppar2} \\
P_{\bot } &=&\sqrt{1-\xi ^{2}}\frac{\partial }{\partial p}-\frac{\sqrt{1-\xi
^{2}}}{p\xi }\left( \xi ^{2}-\frac{k_{\Vert }v_{\Vert }}{\omega }\right) 
\frac{\partial }{\partial \xi }  \label{pperp2}
\end{eqnarray}

After some manipulation, the operators in (\ref{qloporig}) can be expressed
in a divergence form 
\begin{eqnarray}
P_{\Vert } &=&\frac{1}{p^{2}}\frac{\partial }{\partial p}p^{2}\xi -\frac{1}{p%
}\frac{\partial }{\partial \xi }\left( -1+\xi ^{2}+\frac{n\Omega }{\omega }%
\right) \\
P_{\bot }-\frac{1}{p_{\bot }}\left( \frac{k_{\Vert }v_{\Vert }}{\omega }%
-1\right) &=&\frac{1}{p^{2}}\frac{\partial }{\partial p}p^{2}\sqrt{1-\xi ^{2}%
}-\frac{1}{p}\frac{\partial }{\partial \xi }\frac{\sqrt{1-\xi ^{2}}}{\xi }%
\left( \xi ^{2}-\frac{k_{\Vert }v_{\Vert }}{\omega }\right)
\end{eqnarray}
so that (\ref{qloporig}) can be rewritten as 
\begin{equation}
Q\left( f\right) =-\frac{1}{p^{2}}\frac{\partial }{\partial p}\left(
p^{2}S_{p}^{\text{RF}}\right) +\frac{1}{p}\frac{\partial }{\partial \xi }%
\left( \sqrt{1-\xi ^{2}}S_{\xi }^{\text{RF}}\right)  \notag
\end{equation}
with 
\begin{eqnarray}
S_{p}^{\text{RF}} &=&\lim_{V\rightarrow \infty }\frac{e^{2}}{\left( 2\pi
\right) ^{3}}\frac{1}{V}\sum_{n=-\infty }^{\infty }\iiint d^{3}k\left[
\left\{ \xi E_{\mathbf{k},\Vert }^{\ast }J_{n}+\sqrt{1-\xi ^{2}}E_{\mathbf{k}%
,\bot }^{\ast }\right\} \right. \\
&&\left. \cdot \frac{i}{\left[ n\Omega +k_{\Vert }v_{\Vert }-\omega \right] }%
\left\{ E_{\mathbf{k},\Vert }J_{n}P_{\Vert }+E_{\mathbf{k},\bot }P_{\bot
}\right\} f\right]  \notag
\end{eqnarray}
\begin{eqnarray}
S_{\xi }^{\text{RF}} &=&\lim_{V\rightarrow \infty }\frac{e^{2}}{\left( 2\pi
\right) ^{3}}\frac{1}{V}\sum_{n=-\infty }^{\infty }\iiint d^{3}k\left[
\left\{ \frac{-1}{\sqrt{1-\xi ^{2}}}\left( 1-\xi ^{2}-\frac{n\Omega }{\omega
^{\ast }}\right) E_{\mathbf{k},\Vert }^{\ast }J_{n}+\frac{1}{\xi }\left( \xi
^{2}-\frac{k_{\Vert }v_{\Vert }}{\omega ^{\ast }}\right) E_{\mathbf{k},\bot
}^{\ast }\right\} \right. \\
&&\left. \cdot \frac{i}{\left[ n\Omega +k_{\Vert }v_{\Vert }-\omega \right] }%
\left\{ E_{\mathbf{k},\Vert }J_{n}P_{\Vert }+E_{\mathbf{k},\bot }P_{\bot
}\right\} f\right]  \notag
\end{eqnarray}

Considering the operators (\ref{ppar2}-\ref{pperp2}), we wee that the fluxes
in momentum space have a purely diffusive form, and can be expressed in the
form 
\begin{eqnarray}
S_{p}^{\text{RF}} &=&-D_{pp}^{\text{RF}}\frac{\partial f}{\partial p}%
+D_{p\xi }^{\text{RF}}\frac{\sqrt{1-\xi ^{2}}}{p}\frac{\partial f}{\partial
\xi } \\
S_{\xi }^{\text{RF}} &=&-D_{\xi p}^{\text{RF}}\frac{\partial f}{\partial p}%
+D_{\xi \xi }^{\text{RF}}\frac{\sqrt{1-\xi ^{2}}}{p}\frac{\partial f}{%
\partial \xi }
\end{eqnarray}
with 
\begin{eqnarray}
D_{pp}^{\text{RF}} &=&-\lim_{V\rightarrow \infty }\frac{e^{2}}{\left( 2\pi
\right) ^{3}}\frac{1}{V}\sum_{n=-\infty }^{\infty }\iiint d^{3}k\left[
\left\{ \xi E_{\mathbf{k},\Vert }^{\ast }J_{n}+\sqrt{1-\xi ^{2}}E_{\mathbf{k}%
,\bot }^{\ast }\right\} \right. \\
&&\left. \cdot \frac{i}{\left[ n\Omega +k_{\Vert }v_{\Vert }-\omega \right] }%
\left\{ \xi E_{\mathbf{k},\Vert }J_{n}+\sqrt{1-\xi ^{2}}E_{\mathbf{k},\bot
}\right\} \right]  \notag \\
D_{p\xi }^{\text{RF}} &=&\lim_{V\rightarrow \infty }\frac{e^{2}}{\left( 2\pi
\right) ^{3}}\frac{1}{V}\sum_{n=-\infty }^{\infty }\iiint d^{3}k\left[
\left\{ \xi E_{\mathbf{k},\Vert }^{\ast }J_{n}+\sqrt{1-\xi ^{2}}E_{\mathbf{k}%
,\bot }^{\ast }\right\} \right. \\
&&\left. \cdot \frac{i}{\left[ n\Omega +k_{\Vert }v_{\Vert }-\omega \right] }%
\left\{ \frac{1}{\sqrt{1-\xi ^{2}}}\left( 1-\xi ^{2}-\frac{n\Omega }{\omega }%
\right) E_{\mathbf{k},\Vert }J_{n}-\frac{1}{\xi }\left( \xi ^{2}-\frac{%
k_{\Vert }v_{\Vert }}{\omega }\right) E_{\mathbf{k},\bot }\right\} \right] 
\notag \\
D_{\xi p}^{\text{RF}} &=&-\lim_{V\rightarrow \infty }\frac{e^{2}}{\left(
2\pi \right) ^{3}}\frac{1}{V}\sum_{n=-\infty }^{\infty }\iiint d^{3}k\left[
\left\{ \frac{-1}{\sqrt{1-\xi ^{2}}}\left( 1-\xi ^{2}-\frac{n\Omega }{\omega
^{\ast }}\right) E_{\mathbf{k},\Vert }^{\ast }J_{n}+\frac{1}{\xi }\left( \xi
^{2}-\frac{k_{\Vert }v_{\Vert }}{\omega ^{\ast }}\right) E_{\mathbf{k},\bot
}^{\ast }\right\} \right. \\
&&\left. \cdot \frac{i}{\left[ n\Omega +k_{\Vert }v_{\Vert }-\omega \right] }%
\left\{ E_{\mathbf{k},\Vert }J_{n}\xi +E_{\mathbf{k},\bot }\sqrt{1-\xi ^{2}}%
\right\} \right]  \notag \\
D_{\xi \xi }^{\text{RF}} &=&\lim_{V\rightarrow \infty }\frac{e^{2}}{\left(
2\pi \right) ^{3}}\frac{1}{V}\sum_{n=-\infty }^{\infty }\iiint d^{3}k\left[
\left\{ \frac{-1}{\sqrt{1-\xi ^{2}}}\left( 1-\xi ^{2}-\frac{n\Omega }{\omega
^{\ast }}\right) E_{\mathbf{k},\Vert }^{\ast }J_{n}+\frac{1}{\xi }\left( \xi
^{2}-\frac{k_{\Vert }v_{\Vert }}{\omega ^{\ast }}\right) E_{\mathbf{k},\bot
}^{\ast }\right\} \right. \\
&&\left. \cdot \frac{i}{\left[ n\Omega +k_{\Vert }v_{\Vert }-\omega \right] }%
\left\{ \frac{1}{\sqrt{1-\xi ^{2}}}\left( 1-\xi ^{2}-\frac{n\Omega }{\omega }%
\right) E_{\mathbf{k},\Vert }J_{n}-\frac{1}{\xi }\left( \xi ^{2}-\frac{%
k_{\Vert }v_{\Vert }}{\omega }\right) E_{\mathbf{k},\bot }\right\} \right] 
\notag
\end{eqnarray}

\bigskip In the limit of resonant diffusion, we have 
\begin{equation}
\frac{1}{\left[ n\Omega +k_{\Vert }v_{\Vert }-\omega \right] }\rightarrow
i\pi \delta \left( \omega -k_{\Vert }v_{\Vert }-n\Omega \right)
\end{equation}
and we express the resonance condition 
\begin{equation}
k_{\Vert }v_{\Vert }=\omega -n\Omega
\end{equation}
so that we get 
\begin{eqnarray}
D_{pp}^{\text{RF}} &=&\sum_{n=-\infty }^{\infty }\left( 1-\xi ^{2}\right)
D_{n}^{\text{RF}}\left( p,\xi \right)  \notag \\
D_{p\xi }^{\text{RF}} &=&\sum_{n=-\infty }^{\infty }-\frac{\sqrt{1-\xi ^{2}}%
}{\xi }\left( 1-\xi ^{2}-\frac{n\Omega }{\omega }\right) D_{n}^{\text{RF}%
}\left( p,\xi \right) \\
D_{\xi p}^{\text{RF}} &=&\sum_{n=-\infty }^{\infty }-\frac{\sqrt{1-\xi ^{2}}%
}{\xi }\left( 1-\xi ^{2}-\frac{n\Omega }{\omega }\right) D_{n}^{\text{RF}%
}\left( p,\xi \right)  \notag \\
D_{\xi \xi }^{\text{RF}} &=&\sum_{n=-\infty }^{\infty }\frac{1}{\xi ^{2}}%
\left( 1-\xi ^{2}-\frac{n\Omega }{\omega }\right) ^{2}D_{n}^{\text{RF}%
}\left( p,\xi \right)  \notag
\end{eqnarray}
with 
\begin{equation}
D_{n}^{\text{RF}}\left( p,\xi \right) =\lim_{V\rightarrow \infty }\frac{\pi
e^{2}}{V}\iiint \frac{d^{3}k}{\left( 2\pi \right) ^{3}}\delta \left( \omega
-k_{\Vert }v_{\Vert }-n\Omega \right) \left| E_{\mathbf{k},\bot }+\frac{\xi 
}{\sqrt{1-\xi ^{2}}}E_{\mathbf{k},\Vert }J_{n}\right| ^{2}  \label{dnrf}
\end{equation}

Using (\ref{eperp}), we can define 
\begin{equation}
D_{\mathbf{k}}^{(n)}=e^{2}\left| \frac{1}{\sqrt{2}}E_{\mathbf{k}%
,+}e^{-i\alpha }J_{n-1}\left( \frac{k_{\bot }v_{\bot }}{\Omega }\right) +%
\frac{1}{\sqrt{2}}E_{\mathbf{k},-}e^{+i\alpha }J_{n+1}\left( \frac{k_{\bot
}v_{\bot }}{\Omega }\right) +\frac{p_{\Vert }}{p_{\bot }}E_{\mathbf{k},\Vert
}J_{n}\left( \frac{k_{\bot }v_{\bot }}{\Omega }\right) \right| ^{2}
\label{dk}
\end{equation}
so that (\ref{dnrf}) becomes 
\begin{equation}
D_{n}^{\text{RF}}(p,\xi )=\lim_{V\rightarrow \infty }\frac{\pi }{V}\int 
\frac{d^{3}k}{(2\pi )^{3}}D_{\mathbf{k}}^{(n)}\delta (\omega -k_{\Vert
}v_{\Vert }-n\Omega )
\end{equation}

\section{\protect\bigskip QL Diffusion Coefficient for a Plane Wave\label%
{plane}}

We take the limit case of a plane wave for a given ray, with given frequency 
$\omega _{b}$ and wave number $\mathbf{k}_{b}$. If we assume

\begin{eqnarray}
\mathbf{E}_{b}\left( \mathbf{x},t\right) &=&\text{Re}\left[ \mathbf{E}%
_{b0}e^{i\left( \mathbf{k}_{b}\cdot \mathbf{x}-\omega _{b}t\right) }\right]
\\
&=&\frac{1}{2}\left[ \mathbf{E}_{b0}e^{i\left( \mathbf{k}_{b}\cdot \mathbf{x}%
-\omega _{b}t\right) }+\mathbf{E}_{b0}^{\ast }e^{-i\left( \mathbf{k}%
_{b}\cdot \mathbf{x}-\omega _{b}t\right) }\right]
\end{eqnarray}
so that 
\begin{eqnarray}
\mathbf{E}_{b\mathbf{k}}\left( \mathbf{k},t\right) &\equiv &\iiint
d^{3}x\;e^{-i\mathbf{k}\cdot \mathbf{x}}\mathbf{E}_{b}\left( \mathbf{x}%
,t\right)  \notag \\
&=&\frac{1}{2}\left[ \mathbf{E}_{b0}e^{-i\omega _{b}t}\iiint
d^{3}x\;e^{i\left( \mathbf{k}_{b}-\mathbf{k}\right) \cdot \mathbf{x}}+%
\mathbf{E}_{b0}^{\ast }e^{i\omega _{b}t}\iiint d^{3}x\;e^{i\left( \mathbf{k}%
_{b}+\mathbf{k}\right) \cdot \mathbf{x}}\right]  \notag \\
&=&\frac{1}{2}\left[ \mathbf{E}_{b0}e^{-i\omega _{b}t}\delta \left( \mathbf{k%
}_{b}-\mathbf{k}\right) +\mathbf{E}_{b0}^{\ast }e^{i\omega _{b}t}\delta
\left( \mathbf{k}_{b}+\mathbf{k}\right) \right]
\end{eqnarray}

We therefore get 
\begin{equation}
D_{b,n}^{\text{RF}}(p,\xi )=\lim_{V\rightarrow \infty }\frac{q_{e}^{2}\pi }{V%
}\iiint \frac{d^{3}k}{(2\pi )^{3}}\left| \mathbf{E}_{b\mathbf{k}}\right|
^{2}\left| \Theta _{\mathbf{k}}^{b,n}\left( \mathbf{k}\right) \right|
^{2}\delta \left( \omega _{b}-k_{\Vert }v_{\Vert }-n\Omega \right)
\end{equation}

We have 
\begin{eqnarray}
\left| \mathbf{E}_{b\mathbf{k}}\right| ^{2} &=&\frac{1}{4}\left| \mathbf{E}%
_{b0}e^{-i\omega _{b}t}\delta \left( \mathbf{k}_{b}-\mathbf{k}\right) +%
\mathbf{E}_{b0}^{\ast }e^{i\omega _{b}t}\delta \left( \mathbf{k}_{b}+\mathbf{%
k}\right) \right| ^{2}  \notag \\
&=&\frac{1}{4}\left( \mathbf{E}_{b0}e^{-i\omega _{b}t}\delta \left( \mathbf{k%
}_{b}-\mathbf{k}\right) +\mathbf{E}_{b0}^{\ast }e^{i\omega _{b}t}\delta
\left( \mathbf{k}_{b}+\mathbf{k}\right) \right) \left( \mathbf{E}_{b0}^{\ast
}e^{i\omega _{b}t}\delta \left( \mathbf{k}_{b}-\mathbf{k}\right) +\mathbf{E}%
_{b0}e^{-i\omega _{b}t}\delta \left( \mathbf{k}_{b}+\mathbf{k}\right) \right)
\notag \\
&=&\frac{1}{4}\left( \left| \mathbf{E}_{b0}\right| ^{2}\delta ^{2}\left( 
\mathbf{k}_{b}-\mathbf{k}\right) +\left| \mathbf{E}_{b0}\right| ^{2}\delta
^{2}\left( \mathbf{k}_{b}+\mathbf{k}\right) +2 \text{Re}\left[ \mathbf{E}%
_{b0}^{2}e^{-i\omega _{b}t}\right] \delta \left( \mathbf{k}_{b}-\mathbf{k}%
\right) \delta \left( \mathbf{k}_{b}+\mathbf{k}\right) \right)
\end{eqnarray}

Expressing the condition $\mathbf{k}=\pm \mathbf{k}_{b}$ from the first
delta function, we get 
\begin{eqnarray}
D_{b,n}^{\text{RF}}(p,\xi ) &=&\frac{1}{4}q_{e}^{2}\pi \left| \Theta _{%
\mathbf{k}}^{b,n}\left( \mathbf{k}_{b}\right) \right| ^{2}\delta \left(
\omega _{b}-k_{\Vert b}v_{\Vert }-n\Omega \right) \lim_{V\rightarrow \infty }%
\frac{1}{V}\iiint \frac{d^{3}k}{(2\pi )^{3}} \\
&&\left[ \left| \mathbf{E}_{b0}\right| ^{2}\delta ^{2}\left( \mathbf{k}_{b}-%
\mathbf{k}\right) +\left| \mathbf{E}_{b0}\right| ^{2}\delta ^{2}\left( 
\mathbf{k}_{b}+\mathbf{k}\right) +2\text{Re}\left[ \mathbf{E}%
_{b0}^{2}e^{-i\omega _{b}t}\right] \delta \left( \mathbf{k}_{b}-\mathbf{k}%
\right) \delta \left( \mathbf{k}_{b}+\mathbf{k}\right) \right]
\end{eqnarray}

The term $\delta \left( \mathbf{k}_{b}-\mathbf{k}\right) \delta \left( 
\mathbf{k}_{b}+\mathbf{k}\right) $ has no contribution when the limit $%
V\rightarrow \infty $ is taken. Now, using Parseval's theorem, 
\begin{eqnarray}
\lim_{V\rightarrow \infty }\frac{1}{V}\iiint \frac{d^{3}k}{(2\pi )^{3}}%
\left( \delta ^{2}\left( \mathbf{k}_{b}-\mathbf{k}\right) +\delta ^{2}\left( 
\mathbf{k}_{b}+\mathbf{k}\right) \right) &=&2\lim_{V\rightarrow \infty }%
\frac{1}{V}\iiint d^{3}x  \notag \\
&=&2
\end{eqnarray}

We finally get 
\begin{equation}
D_{b,n}^{\text{RF}}(p,\xi )=q_{e}^{2}\pi \frac{\left| \mathbf{E}_{b0}\right|
^{2}}{2}\left| \Theta _{\mathbf{k}}^{b,n}\left( \mathbf{k}_{b}\right)
\right| ^{2}\delta \left( \omega _{b}-k_{\Vert b}v_{\Vert }-n\Omega \right)
\end{equation}

\section{Cold Plasma Description for RF Waves\label{cold}}

\subsection{Introduction}

In the cold plasma description, a fluid-like response of the plasma to the
field is assumed. This description requires that

\begin{itemize}
\item  We remain far from resonances, where $k_{\Vert }\rightarrow \infty $.
Near resonances, mode conversion to electrostatic waves (IBW, EBW) occur,
which are not described by cold plasma wave theory. LHCD and ECCD usually
take place far enough from the LH and UH resonances, respectively, so that
the cold plasma description is valid.

\item  Also, when $\omega /k_{\Vert }\sim v_{Te},v_{Ti}$, strong dissipation
accurs trough Landau Damping and must be considered. In LHCD, the
wave-particle interaction usually takes place for $\omega /k_{\Vert
}>3v_{Te} $ and therefore cold plasma description still holds for LHW.

\item  FLR effects must be small, that is, we must have 
\begin{equation}
\left| z\right| =\left| \frac{k_{\perp }v_{\perp }}{\Omega }\right| \ll 1
\end{equation}

As we saw in previous sections, this condition usually holds for LHCD, given
that the temperature is not too high, and also for ECCD, if in addition the
harmonic is low ($n=1$ or $2$).
\end{itemize}

\subsection{Power Flux}

The relation between power flux and electric field was described using the
vector $\mathbf{\Phi }_{b}$ defined in (\ref{phib}) 
\begin{equation}
\mathbf{\Phi }_{b}=\mathbf{\Phi }_{bP}+\mathbf{\Phi }_{bT}
\end{equation}
with 
\begin{equation}
\mathbf{\Phi }_{bP}=\frac{\mu _{0}c}{\left| \mathbf{E}_{b0}\right| ^{2}}%
\text{Re}\left[ \widetilde{\mathbf{E}}_{b}\times \widetilde{\mathbf{H}}%
_{b}^{\ast }\right]  \label{phibp}
\end{equation}
and 
\begin{equation}
\mathbf{\Phi }_{bT}=\mathbf{-}\frac{1}{2}\mathbf{e}_{b}^{\ast }\cdot \frac{%
\partial \mathbb{K}^{H}}{\partial \mathbf{N}}\cdot \mathbf{e}_{b}
\end{equation}

In the cold plasma model, the dielectric tensor $\mathbb{K}$ is independent
of $\mathbf{N}$ and the contribution $\mathbf{\Phi }_{bT}$, from the
coherent motion of particles, vanishes. It is usually a very good
approximation to neglect this kinetic power flux, even for
quasi-electrostatic waves like Lower-Hybrid waves (See Appendix \ref{kinetic}%
). 
\begin{equation}
\mathbf{\Phi }_{bT}=0
\end{equation}

The complex electric field is 
\begin{equation}
\widetilde{\mathbf{E}}_{b}\left( x,t\right) =\mathbf{E}_{b0}e^{i\mathbf{k}%
_{b}\cdot \mathbf{x-}\omega _{b}t}
\end{equation}
and the magnetic field is given by 
\begin{equation}
\widetilde{\mathbf{H}}_{b}^{\ast }\mathbf{=}\frac{\mathbf{k}_{b}\times 
\widetilde{\mathbf{E}}_{b}}{\mu _{0}\omega _{b}}=\frac{\mathbf{N}_{b}\times 
\widetilde{\mathbf{E}}_{b}}{\mu _{0}c}
\end{equation}
so that (\ref{phibp}) becomes 
\begin{eqnarray}
\mathbf{\Phi }_{bP} &=&\text{Re}\left[ \mathbf{e}_{b}\times \left( \mathbf{N}%
_{b}\times \mathbf{e}_{b}^{\ast }\right) \right]  \notag \\
&=&\text{Re}\left[ \left| e_{b}\right| ^{2}\mathbf{N}_{b}-\left( \mathbf{N}%
_{b}\cdot \mathbf{e}_{b}\right) \mathbf{e}_{b}^{\ast }\right]  \label{phibp2}
\end{eqnarray}
where $\mathbf{\varepsilon }$ is the polarization vector, which can be
decomposed into its components parallel and perpendicular to $\mathbf{N}_{b}$%
: 
\begin{eqnarray}
\mathbf{e}_{bL} &=&\left( \mathbf{e}_{b}\cdot \widehat{N}_{b}\right) 
\widehat{N}_{b} \\
\mathbf{e}_{bT} &=&\mathbf{e}_{b}-\mathbf{e}_{bL}
\end{eqnarray}

Then (\ref{phibp2}) becomes 
\begin{eqnarray}
\mathbf{\Phi }_{bP} &=&\text{Re}\left[ \left( \left| e_{bL}\right|
^{2}+\left| e_{bT}\right| ^{2}\right) \mathbf{N}_{b}-Ne_{bL}\left( \mathbf{e}%
_{bL}^{\ast }+\mathbf{e}_{bT}^{\ast }\right) \right]  \notag \\
&=&\text{Re}\left[ \left| e_{bT}\right| ^{2}\mathbf{N}_{b}-N_{b}e_{bL}%
\mathbf{e}_{bT}^{\ast }\right]  \notag \\
&=&\left| e_{bT}\right| ^{2}\mathbf{N}_{b}-N_{b}\text{Re}\left( e_{bL}%
\mathbf{e}_{bT}^{\ast }\right)  \label{phibpres}
\end{eqnarray}
so that finally 
\begin{equation}
\mathbf{\Phi }_{b}=\left| e_{bT}\right| ^{2}\mathbf{N}_{b}-N_{b}\text{Re}%
\left( e_{bL}\mathbf{e}_{bT}^{\ast }\right)  \label{phibres}
\end{equation}

\subsection{Wave Equation and Dispersion Tensor}

We consider the two-fluids description of a non relativistic plasma in a
constant magnetic field $\mathbf{B}_{0}=B_{0}\widehat{z}$. In a continuous
homogeneous linear medium, a Fourier component $\mathbf{E}\left( \mathbf{k}%
,\omega \right) $ verifies the wave equation \cite{ber01c6}

\begin{equation}
\mathbf{N}\times (\mathbf{N}\times \mathbf{E})+\mathbb{K}\cdot \mathbf{E}=0
\end{equation}
or equivalently 
\begin{equation}
(\mathbf{N}\cdot \mathbf{E})\mathbf{N}-N^{2}\mathbf{E}+\mathbb{K}\cdot 
\mathbf{E}=0
\end{equation}
where $\mathbf{N}=c\mathbf{k}/\omega $ is the wave refractive index and $%
\mathbb{K}$ is the dielectric tensor. This equation can be expressed in
tensorial form as 
\begin{equation}
\mathbb{D}\cdot \mathbf{E}=0  \label{waveeq3}
\end{equation}
where the dispersion tensor is 
\begin{equation}
\mathbb{D=}\mathbf{NN-}N^{2}\mathbb{I}+\mathbb{K}
\end{equation}

By evaluating the susceptibilities in the cold plasma limit, the following
dielectric tensor is obtained \cite{bon84} in the frame $\left( x,y,z\right) 
$: 
\begin{equation}
\mathbb{K}= 
\begin{pmatrix}
K_{\perp } & -iK_{xy} & 0 \\ 
iK_{xy} & K_{\perp } & 0 \\ 
0 & 0 & K_{\Vert }
\end{pmatrix}
\label{ktensor}
\end{equation}
whose elements in a plasma made of species $s$ are given by 
\begin{eqnarray}
K_{\perp } &=&1-\sum_{s}\frac{\omega _{ps}^{2}}{\omega ^{2}-\Omega _{s}^{2}}
\notag \\
K_{\parallel } &=&1-\sum_{s}\frac{\omega _{ps}^{2}}{\omega ^{2}}  \notag \\
K_{xy} &=&\sum_{s}\frac{\omega _{ps}^{2}\Omega }{\omega (\omega ^{2}-\Omega
_{s}^{2})}
\end{eqnarray}
where 
\begin{equation}
\omega _{ps}=\sqrt{\frac{n_{s}q_{s}^{2}}{\varepsilon _{0}m_{s}}}
\end{equation}
is the plasma frequency and 
\begin{equation}
\Omega _{s}=\frac{q_{s}B_{0}}{m_{s}}
\end{equation}
the cyclotron frequency for the species $s$.

\subsection{Dispersion Relation}

In order to have a non-trivial solution to the wave equation, the dispersion
relation must therefore be satisfied: 
\begin{equation}
D\left( \mathbf{k,}\omega \right) =\left| \mathbb{D}\right| =0
\label{disprel}
\end{equation}

The axis $\widehat{x}$ is chosen such that 
\begin{equation}
\mathbf{N=}N_{\perp }\widehat{x}+N_{\parallel }\widehat{z}
\end{equation}

Taking the parallel component $N_{\parallel }$ as given, (\ref{disprel})
lead to the following equation for $N_{\perp }$%
\begin{equation}
AN_{\perp }^{4}+BN_{\perp }^{2}+C=0
\end{equation}
with 
\begin{eqnarray}
A &=&K_{\perp } \\
B &=&\left( N_{\parallel }^{2}-K_{\perp }\right) \left( K_{\perp
}+K_{\parallel }\right) +K_{xy}^{2} \\
C &=&K_{\parallel }\left[ \left( N_{\parallel }^{2}-K_{\perp }\right)
^{2}-K_{xy}^{2}\right]
\end{eqnarray}
or equivalently 
\begin{eqnarray}
A &=&K_{\perp } \\
B &=&\left( K_{\perp }+K_{\parallel }\right) N_{\parallel }^{2}-\left(
K_{R}K_{L}+K_{\perp }K_{\parallel }\right) \\
C &=&K_{\parallel }\left( N_{\parallel }^{2}-K_{R}\right) \left(
N_{\parallel }^{2}-K_{L}\right)
\end{eqnarray}
with 
\begin{eqnarray}
K_{R} &=&K_{\perp }+K_{xy} \\
K_{L} &=&K_{\perp }-K_{xy}
\end{eqnarray}

We find the expression for $N_{\perp }$ as a function of $N_{\parallel }$
and $\omega $: 
\begin{equation}
N_{\perp }^{2}=\frac{-B\pm \sqrt{B^{2}-4AC}}{2A}
\end{equation}

We see immediately that the resonances occur for $A=0$, that is, 
\begin{equation}
K_{\perp }=0
\end{equation}
and that cut-off occur for $C=0$, that is, 
\begin{equation}
K_{R}=N_{\parallel }^{2}
\end{equation}
\begin{equation}
K_{L}=N_{\parallel }^{2}
\end{equation}
or 
\begin{equation}
K_{\parallel }=0
\end{equation}

\subsection{Polarization components}

Once $N_{\perp }$ has been evaluated, the components of the polarization of
the electric field are simply given as the eigenvector of the wave equation (%
\ref{waveeq3}) 
\begin{equation}
\mathbb{D}\cdot \mathbf{E}=0
\end{equation}

\section{Electrostatic Approximation for LH Waves\label{lhw}}

\subsection{Electrostatic Dispersion Relation}

We consider the wave equation (\ref{waveeq3}) 
\begin{equation}
(\mathbf{N}\cdot \mathbf{E})\mathbf{N}-N^{2}\mathbf{E}+\mathbb{K}\cdot 
\mathbf{E}=0  \label{waveeq}
\end{equation}

The electric field can be separated into its longitudinal and transverse
components with respect to the normalized wave vector $\mathbf{N}$: 
\begin{equation}
\mathbf{E}=\mathbf{E}_{L}+\mathbf{E}_{T}  \label{edecomp}
\end{equation}
so that the wave equation (\ref{waveeq}) becomes 
\begin{equation}
(N^{2}-\mathbb{K})\mathbf{E}_{T}-\mathbb{K}\cdot \mathbf{E}_{L}=0
\label{waveeq2}
\end{equation}

Electrostatic waves, and also electromagnetic waves approaching resonances,
can satisfy the condition 
\begin{equation}
N^{2}\gg |K_{ij}|  \label{esapprox}
\end{equation}
and therefore (\ref{waveeq2}) reduces to 
\begin{equation}
N^{2}\mathbf{E}_{T}-\mathbb{K}\cdot \mathbf{E}_{L}\simeq 0  \label{waveeqes}
\end{equation}

Dot-multiplying (\ref{waveeqes}) by $\widehat{N}=\mathbf{N}/N$ leads to the
electrostatic dispersion relation 
\begin{equation}
D_{L}\equiv \widehat{N}\cdot \mathbb{K}\cdot \widehat{N}=0  \label{dispreles}
\end{equation}
and the transverse electric field is given by 
\begin{equation}
\mathbf{E}_{T}=\frac{1}{N^{2}}\mathbb{K}\cdot \mathbf{E}_{L}  \label{et}
\end{equation}

We note from (\ref{et}) and (\ref{esapprox}) that $|\mathbf{E}_{T}|\ll |%
\mathbf{E}_{L}|$. The electric field is quasi longitudinal, which justify
the term electrostatic approximation given to the condition (\ref{esapprox}).

\subsection{Cold Plasma Limit}

Using (\ref{ktensor}), the electrostatic dispersion relation (\ref{dispreles}%
) in the cold plasma limit is then given by 
\begin{equation}
D_{L}=N_{\perp }^{2}K_{\perp }+N_{\parallel }^{2}K_{\parallel }=0
\label{disprelescpl}
\end{equation}
which gives an expression for $N_{\perp }$ as a function of $N_{\parallel }$
and $\omega $: 
\begin{equation}
N_{\perp }^{2}=\frac{-K_{\parallel }}{K_{\perp }}N_{\parallel }^{2}
\label{nperplh0}
\end{equation}

Since the electric field is quasi-longitudinal ($\mathbf{E}\parallel \mathbf{%
N}$), we also have the following relation for the polarization vector $%
\mathbf{e}$ 
\begin{equation}
\frac{\left| e_{\parallel }\right| }{\left| e_{\perp }\right| }=\frac{\left|
N_{\parallel }\right| }{\left| N_{\perp }\right| }  \label{eratio}
\end{equation}

\subsection{Lower Hybrid Waves}

We consider the lower-hybrid range of frequency in a electron-ion plasma ,
where the following ordering applies 
\begin{equation}
\omega _{ci}\ll \omega \ll \omega _{ce}
\end{equation}

Assuming in addition that 
\begin{equation}
\omega \ll \omega _{pe}
\end{equation}
leads to approximate expressions of the dielectric tensor components given
by: 
\begin{eqnarray}
K_{\perp } &\simeq &1+\frac{\omega _{pe}^{2}}{\omega _{ce}^{2}}-\frac{\omega
_{pi}^{2}}{\omega ^{2}}  \notag \\
K_{\parallel } &\simeq &-\frac{\omega _{pe}^{2}}{\omega ^{2}}  \notag \\
K_{xy} &\simeq &\frac{\omega _{pe}^{2}}{\omega \omega _{ce}}
\label{disptenlh}
\end{eqnarray}
The perpendicular index of refraction becomes 
\begin{equation}
N_{\perp }^{2}=\frac{\omega _{pe}^{2}}{\omega ^{2}}\left( 1+\frac{\omega
_{pe}^{2}}{\omega _{ce}^{2}}-\frac{\omega _{pi}^{2}}{\omega ^{2}}\right)
^{-1}N_{\parallel }^{2}
\end{equation}
which can be rewritten as 
\begin{eqnarray}
N_{\perp }^{2} &=&\frac{\omega _{pe}^{2}/\omega ^{2}}{\left( 1+\omega
_{pe}^{2}/\omega _{ce}^{2}\right) \left( 1-\omega _{\text{LH}}^{2}/\omega
^{2}\right) }N_{\parallel }^{2} \\
&=&\frac{\omega _{pe}^{2}/\omega _{pi}^{2}}{\left( \omega ^{2}/\omega _{%
\text{LH}}^{2}-1\right) }N_{\parallel }^{2}  \label{nperplh2}
\end{eqnarray}
where 
\begin{equation}
\omega _{\text{LH}}^{2}=\frac{\omega _{pi}^{2}}{1+\omega _{pe}^{2}/\omega
_{ce}^{2}}
\end{equation}
is the lower hybrid frequency. We recognize that $N_{\perp }\rightarrow
\infty $ as the the lower-hybrid frequency approaches the wave frequency.

\section{Kinetic Power Flux for LHW \label{kinetic}}

\bigskip In the analysis above the kinetic part of the power flux, due to
the coherent motion of charge carriers, has been neglected. This
approximation must be justified by comparing the Poynting and the kinetic
fluxes.

The normalized expression for the kinetic power flux associated with an
electromagnetic wave in a kinetic plasma is given by (\ref{phit}) 
\begin{equation}
\mathbf{\Phi }_{bT}=\mathbf{-}\frac{\omega }{2c}\mathbf{\varepsilon }%
_{b}^{\ast }\cdot \frac{\partial \mathbb{K}^{H}}{\partial \mathbf{k}}\cdot 
\mathbf{\varepsilon }_{b}  \label{kinflux}
\end{equation}

\subsection{Electrostatic Waves}

In the electrostatic approximation, the electric field is quasi-longitudinal
(\ref{et}) 
\begin{equation}
\left| \mathbf{\varepsilon }_{bT}\right| \ll \left| \varepsilon _{bL}\right|
\simeq 1
\end{equation}
and (\ref{kinflux}) becomes 
\begin{equation}
\mathbf{\Phi }_{bT}=\mathbf{-}\frac{\omega }{2c}\widehat{k}\cdot \frac{%
\partial \mathbb{K}^{H}}{\partial \mathbf{k}}\cdot \widehat{k}
\label{kinfluxes1}
\end{equation}

In the frame of the wave vector $\vec{k}$, (\ref{kinfluxes1}) can then be
expressed as 
\begin{equation}
\mathbf{\Phi }_{bT}=\mathbf{-}\frac{\omega }{2c}\frac{\partial K_{L}^{H}}{%
\partial k}\widehat{k}  \label{kinfluxes2}
\end{equation}
where 
\begin{equation}
K_{L}=\widehat{k}\cdot \mathbb{K}\cdot \widehat{k}
\end{equation}
is the longitudinal (electrostatic) component of the dielectric tensor.

\subsection{Lower Hybrid Waves}

An expression of the electrostatic dispersion relation for LH waves that
includes the first-order thermal corrections is given by \cite{por02}: 
\begin{equation}
K_{L}^{\text{LH}}(\omega ,k)\simeq 1+\frac{k_{\perp }^{2}}{k^{2}}\frac{%
\omega _{pe}^{2}}{\Omega _{e}^{2}}\left( 1-\frac{3b_{e}}{4}\right) -\frac{%
k_{\perp }^{2}}{k^{2}}\frac{\omega _{pi}^{2}}{\omega ^{2}}\left( 1+\frac{%
3b_{i}\Omega _{i}^{2}}{\omega ^{2}}\right) -\frac{k_{\parallel }^{2}}{k^{2}}%
\frac{\omega _{pe}^{2}}{\omega ^{2}}  \label{dispreleslhkin}
\end{equation}
where the finite Larmor radius (FRL) effects are scaled by the parameters 
\begin{equation}
b_{e}=\frac{k_{\perp }^{2}v_{Te}^{2}}{\Omega _{e}^{2}},\qquad b_{i}=\frac{%
k_{\perp }^{2}v_{Ti}^{2}}{\Omega _{i}^{2}}  \label{flr}
\end{equation}
with $v_{T}^{2}=T/m$.

The $k$-dependence of the dispersion relation is found in the thermal
correction terms so that 
\begin{equation}
\frac{\partial K_{L}^{\text{LH}}}{\partial k}=-\frac{3}{4}\frac{k_{\perp
}^{4}}{k^{4}}\frac{\omega _{pe}^{2}}{\Omega _{e}^{2}}\frac{2kv_{Te}^{2}}{%
\Omega _{e}^{2}}-\frac{k_{\perp }^{4}}{k^{4}}\frac{\omega _{pi}^{2}}{\omega
^{2}}\frac{3\Omega _{i}^{2}}{\omega ^{2}}\frac{2kv_{Ti}^{2}}{\Omega _{i}^{2}}
\label{kdepdispreleslhkin}
\end{equation}

Inside the plasma, we have $k_{\parallel }^{2}\ll k_{\perp }^{2}$ and
therefore $k_{\perp }^{2}/k^{2}\simeq 1$. We finally get the following
expression for the kinetic flux associated with quasi-electrostatic LH
waves: 
\begin{equation}
\mathbf{\Phi }_{bT}^{\text{LH}}=\frac{1}{N}\left( \frac{1}{4}\frac{\omega
_{pe}^{2}k^{2}v_{Te}^{2}}{\Omega _{e}^{4}}+\frac{\omega
_{pi}^{2}k^{2}v_{Ti}^{2}}{\omega ^{4}}\right) \widehat{k}
\label{kinfluxeslh}
\end{equation}

The kinetic flux is oriented in the direction of the wave vector, and the
incident flux on the flux surface 
\begin{equation}
\mathbf{\Phi }_{bT\perp }^{\text{LH}}=\frac{1}{N}\left( \frac{1}{4}\frac{%
\omega _{pe}^{2}k^{2}v_{Te}^{2}}{\Omega _{e}^{4}}+\frac{\omega
_{pi}^{2}k^{2}v_{Ti}^{2}}{\omega ^{4}}\right) \left| \widehat{k}\cdot 
\widehat{n}\right|  \label{kinfluxeslhperp}
\end{equation}
This incident kinetic flux must be compared to the incident Poynting flux
taken in the cold plasma limit from (\ref{phiblh2}) 
\begin{equation}
\mathbf{\Phi }_{bP\perp }^{\text{LH}}=\frac{K_{\perp }}{N^{2}}\left| \mathbf{%
N}_{\perp }\cdot \widehat{n}\right| \simeq \frac{K_{\perp }}{N}\left| 
\widehat{k}\cdot \widehat{n}\right|  \label{poyntfluxeslhperp}
\end{equation}

The ratio of the kinetic power flux to the Poynting flux is given by 
\begin{equation}
\frac{\mathbf{\Phi }_{bT\perp }^{\text{LH}}}{\mathbf{\Phi }_{bP\perp }^{%
\text{LH}}}=\frac{1}{K_{\perp }}\left( \frac{3}{4}\frac{\omega
_{pe}^{2}v_{Te}^{2}k^{2}}{\Omega _{e}^{4}}+3\frac{\omega
_{pi}^{2}v_{Ti}^{2}k^{2}}{\omega ^{4}}\right)  \label{fluxratio}
\end{equation}

Using ({\ref{nperplh0}) we have 
\begin{equation}
k^{2}\simeq \frac{\omega ^{2}}{c^{2}}N_{\perp }^{2}\simeq \frac{\omega ^{2}}{%
c^{2}}\frac{\left( -K_{\parallel }\right) }{K_{\perp }}N_{\parallel
}^{2}\simeq \frac{1}{K_{\perp }}\frac{\omega _{pe}^{2}}{c^{2}}N_{\parallel
}^{2}
\end{equation}
so that 
\begin{equation}
\frac{\mathbf{\Phi }_{bT\perp }^{\text{LH}}}{\mathbf{\Phi }_{bP\perp }^{%
\text{LH}}}=\frac{3\beta _{Te}^{2}N_{\parallel }^{2}}{(1+\omega
_{pe}^{2}/\Omega _{e}^{2})(1-\omega _{\text{LH}}^{2}/\omega ^{2})}\left( 
\frac{1}{4}\frac{\omega _{pe}^{4}}{\Omega _{e}^{4}}+\frac{\omega _{pi}^{4}}{%
\omega ^{4}}\frac{T_{i}}{T_{e}}\right)  \label{fluxratio2}
\end{equation}
where $\beta _{Te}^{2}=T_{e}/mc^{2}$. }

It can be seen from (\ref{fluxratio2}) that the kinetic part of the power
flux becomes significant only as the wave approaches the lower-hybrid
resonance very closely. In a typical LHCD context (for instance Alcator
C-Mod, \cite{bon00}), we have 
\begin{equation}
\omega \sim 2\omega _{\text{LH}},\ \omega _{pe}\sim \Omega _{e},\ \omega
_{pi}\sim \omega ,\ T_{e}\sim T_{i}\text{ and }\beta _{Te}\sim 0.1
\label{cmod}
\end{equation}
so that the kinetic part of the power flux is not more than a few percents.

\bibliographystyle{unsrt}
\bibliography{c:/users/joan/papers/personal/bibliography/lhcdjd,c:/users/joan/papers/personal/bibliography/eccdjd,c:/users/joan/papers/personal/bibliography/fpqljd,c:/users/joan/papers/personal/bibliography/hxrsjd,c:/users/joan/papers/personal/bibliography/neocjd,c:/users/joan/papers/personal/bibliography/tranjd,c:/users/joan/papers/personal/bibliography/miscjd,c:/users/joan/papers/personal/bibliography/rawsjd,c:/users/joan/papers/personal/bibliography/rfbcjd}

\end{document}
